\documentclass[12pt,a4paper]{article}

% ── Packages ──────────────────────────────────────────────────────────────────
\usepackage[T1]{fontenc}
\usepackage[utf8]{inputenc}
\usepackage[english]{babel}
\usepackage{lmodern}
\usepackage{geometry}
\usepackage{graphicx}
\usepackage{xcolor}
\usepackage{hyperref}
\usepackage{listings}
\usepackage{booktabs}
\usepackage{longtable}
\usepackage{array}
\usepackage{enumitem}
\usepackage{titlesec}
\usepackage{fancyhdr}
\usepackage{parskip}
\usepackage{microtype}
\usepackage{caption}
\usepackage{float}
\usepackage{amsmath}
\usepackage{mdframed}
\usepackage{tabularx}

% ── Page geometry ─────────────────────────────────────────────────────────────
\geometry{
  top=2.5cm, bottom=2.5cm,
  left=2.5cm, right=2.5cm,
  headheight=15pt
}

% ── Colours ───────────────────────────────────────────────────────────────────
\definecolor{f1red}{HTML}{E10600}
\definecolor{codebg}{HTML}{1F1F27}
\definecolor{codetext}{HTML}{E8E8F0}
\definecolor{codekw}{HTML}{FF7B72}
\definecolor{codestr}{HTML}{A5D6FF}
\definecolor{codecomment}{HTML}{8B949E}
\definecolor{codenumber}{HTML}{79C0FF}
\definecolor{lightgray}{HTML}{F5F5F5}
\definecolor{darkgray}{HTML}{555555}
\definecolor{turtlebg}{HTML}{F0F8F0}
\definecolor{turtlekw}{HTML}{005B00}

% ── Hyperref setup ────────────────────────────────────────────────────────────
\hypersetup{
  colorlinks=true,
  linkcolor=darkgray,
  urlcolor=darkgray,
  citecolor=darkgray,
  pdftitle={F1 Knowledge System -- TP2 Web Semantics Report},
  pdfauthor={Rodrigo Abreu},
}

% ── Section styling ───────────────────────────────────────────────────────────
\titleformat{\section}
  {\large\bfseries}
  {\thesection.}{0.6em}{}
  [\vspace{-0.4em}\color{darkgray}\rule{\linewidth}{1.5pt}]

\titleformat{\subsection}
  {\normalsize\bfseries\color{darkgray}}
  {\thesubsection.}{0.5em}{}

\titleformat{\subsubsection}
  {\normalsize\itshape\color{darkgray}}
  {\thesubsubsection.}{0.5em}{}

% ── Code listing styles ───────────────────────────────────────────────────────
\lstdefinelanguage{SPARQL}{
  morekeywords={PREFIX,SELECT,WHERE,OPTIONAL,FILTER,GROUP,BY,ORDER,LIMIT,
                OFFSET,DISTINCT,COUNT,MIN,MAX,SUM,SAMPLE,INSERT,DELETE,DATA,
                ASC,DESC,UNION,BOUND,STR,CONTAINS,STRSTARTS,IF,IN,AS,FROM,
                GRAPH,CONSTRUCT,ASK,DESCRIBE,HAVING,VALUES,BIND,NOT,EXISTS,
                LCASE,REPLACE,IRI,CONCAT,LANG,LANGMATCHES},
  sensitive=true,
  morestring=[b]",
  morestring=[b]',
  morecomment=[l]{\#},
}

\lstdefinelanguage{Turtle}{
  morekeywords={a,rdf,rdfs,owl,xsd,PREFIX,BASE},
  sensitive=false,
  morestring=[b]",
  morestring=[b]',
  morecomment=[l]{\#},
}

\lstdefinestyle{sparql}{
  language=SPARQL,
  backgroundcolor=\color{codebg},
  basicstyle=\ttfamily\footnotesize\color{codetext},
  keywordstyle=\color{codekw}\bfseries,
  stringstyle=\color{codestr},
  commentstyle=\color{codecomment}\itshape,
  breaklines=true,
  breakatwhitespace=true,
  tabsize=2,
  showstringspaces=false,
  frame=single, framesep=6pt, framerule=0pt, rulecolor=\color{codebg},
  xleftmargin=6pt, xrightmargin=6pt, captionpos=b,
}

\lstdefinestyle{turtle}{
  language=Turtle,
  backgroundcolor=\color{turtlebg},
  basicstyle=\ttfamily\footnotesize\color{black},
  keywordstyle=\color{turtlekw}\bfseries,
  stringstyle=\color{codekw},
  commentstyle=\color{codecomment}\itshape,
  breaklines=true, tabsize=2, showstringspaces=false,
  frame=single, framesep=6pt, framerule=0pt, rulecolor=\color{turtlebg},
  xleftmargin=6pt, xrightmargin=6pt, captionpos=b,
}

\lstdefinestyle{html}{
  backgroundcolor=\color{codebg},
  basicstyle=\ttfamily\footnotesize\color{codetext},
  keywordstyle=\color{codekw}\bfseries,
  stringstyle=\color{codestr},
  breaklines=true, tabsize=2, showstringspaces=false,
  frame=single, framesep=6pt, framerule=0pt, rulecolor=\color{codebg},
  xleftmargin=6pt, xrightmargin=6pt, captionpos=b,
}

\lstdefinestyle{python}{
  language=Python,
  backgroundcolor=\color{codebg},
  basicstyle=\ttfamily\footnotesize\color{codetext},
  keywordstyle=\color{codekw}\bfseries,
  stringstyle=\color{codestr},
  commentstyle=\color{codecomment}\itshape,
  breaklines=true, tabsize=2, showstringspaces=false,
  frame=single, framesep=6pt, framerule=0pt, rulecolor=\color{codebg},
  xleftmargin=6pt, xrightmargin=6pt, captionpos=b,
}

\lstdefinestyle{shell}{
  backgroundcolor=\color{codebg},
  basicstyle=\ttfamily\footnotesize\color{codetext},
  breaklines=true, frame=single, framesep=6pt, framerule=0pt,
  rulecolor=\color{codebg}, xleftmargin=6pt, xrightmargin=6pt,
}

% ── Header / Footer ───────────────────────────────────────────────────────────
\pagestyle{fancy}
\fancyhf{}
\fancyhead[L]{\small\color{darkgray} F1 Knowledge System}
\fancyhead[R]{\small\color{darkgray} Web Semantics --- TP2}
\fancyfoot[C]{\small\thepage}
\renewcommand{\headrulewidth}{0.4pt}
\renewcommand{\footrulewidth}{0pt}
\renewcommand{\headrule}{\color{darkgray}\rule{\linewidth}{0.4pt}}

% ── Info box ──────────────────────────────────────────────────────────────────
\newmdenv[
  backgroundcolor=lightgray,
  linecolor=f1red,
  linewidth=3pt,
  leftline=true, rightline=false, topline=false, bottomline=false,
  innerleftmargin=10pt, innertopmargin=8pt, innerbottommargin=8pt,
]{infobox}

% ─────────────────────────────────────────────────────────────────────────────
\begin{document}
% ─────────────────────────────────────────────────────────────────────────────

% ── Title page ────────────────────────────────────────────────────────────────
\begin{titlepage}
  \centering
  \vspace*{1.5cm}

  {\fontsize{48}{52}\selectfont\bfseries\color{f1red} F\textcolor{black}{1}}\\[0.4cm]
  {\Large\bfseries Knowledge System}\\[0.3cm]
  {\large\color{darkgray} Web Semantics --- Practical Assignment 2 --- Report}

  \vspace{1.5cm}
  \color{darkgray}\rule{0.4\linewidth}{2pt}\\[0.1cm]
  \color{darkgray}\rule{\linewidth}{0.4pt}
  \vspace{1cm}

  \begin{tabular}{rl}
    \textbf{Authors}    & Rodrigo Abreu (113626), Hugo Ribeiro (113402) \\[0.3em]
    \textbf{Course}     & Web Semantics \\[0.3em]
    \textbf{Year}       & 2025--2026 \\[0.3em]
    \textbf{Technology} & Django 6 $\cdot$ GraphDB $\cdot$ OWL 2 DL $\cdot$ SPIN $\cdot$ RDFa $\cdot$ MF2 \\
  \end{tabular}

  \vspace{1.5cm}
  \color{darkgray}\rule{\linewidth}{0.4pt}\\[0.1cm]
  \color{darkgray}\rule{0.4\linewidth}{2pt}

  \vfill
  {\small\color{darkgray}\today}
\end{titlepage}

\tableofcontents
\newpage

% =============================================================================
\section{Introduction}
% =============================================================================

This report covers \textbf{Practical Assignment 2} (TP2), which extends the F1 Knowledge Graph Explorer developed in TP1. The first assignment established the foundation --- a Django application backed by GraphDB with all domain data stored as RDF triples and queried through SPARQL --- but without any formal ontology or reasoning capabilities. TP2 adds a complete semantic knowledge layer on top of that foundation, enabling entity classification, materialisation of derived facts, semantic annotation of HTML pages, and live integration with external linked data sources.

The key contributions of TP2 are:

\begin{itemize}[leftmargin=1.5em]
  \item An \textbf{OWL 2 DL ontology} that formally describes the F1 domain,
        enabling automatic entity classification directly from the data properties
        of each entity --- without any explicit \texttt{rdf:type} assertions in
        the facts file.
  \item \textbf{21 SPIN rules} (SPARQL INSERT WHERE) that materialise
        inferences OWL cannot express: World Champions, Veterans, Historic
        Circuits, podium decomposition (1st/2nd/3rd), pole positions,
        hat tricks, and the teammate relation between drivers who raced together.
  \item \textbf{RDF reification} annotating every championship-winning triple
        with the driver's final points total.
  \item \textbf{Live semantic enrichment} from Wikidata and DBpedia via
        SPARQLWrapper: driver death dates, team logos, circuit lap records, and more.
  \item \textbf{RDFa 1.1 Lite} (Schema.org) and \textbf{Microformats 2} embedded
        in every detail page, making the site machine-readable by search engines
        and IndieWeb tools alike.
  \item New application pages driven entirely by inferred knowledge:
        World Champions, Multi-Champions, Veterans, Constructor Champions,
        Pole Position Leaders, Hat Trick Hall of Fame, and a Hall of Fame dashboard.
\end{itemize}

\subsection{Technology Stack}

\begin{center}
\begin{tabular}{@{}lll@{}}
\toprule
\textbf{Technology} & \textbf{Role} & \textbf{Where used} \\
\midrule
Python 3.12 / Django 5    & Web framework and data pipeline            & All views, management commands \\
GraphDB 10 (Ontotext)     & RDF triple store with OWL-Max reasoner     & Primary knowledge base \\
SPARQL 1.1                & Knowledge retrieval and SPIN rules         & 21 rules; all Django views \\
RDF / N-Triples           & Data serialisation                         & \texttt{formula1\_integrated.nt} (6.6M triples) \\
RDFS                      & Schema and subclass hierarchy              & Ontology: domain, range, subClassOf \\
OWL 2 DL                  & Ontology axioms and offline reasoning      & \texttt{formula1\_ontology.ttl} (20 classes) \\
Protégé + HermiT 1.4.3    & Ontology editing and consistency check     & Offline validation \\
SPIN (via SPARQL UPDATE)  & Materialised inference beyond OWL          & \texttt{spin\_rules.py} (21 rules) \\
SPARQLWrapper 2           & Federated queries to Wikidata / DBpedia    & \texttt{services/external\_data.py} \\
RDFa 1.1 Lite             & Semantic HTML annotation (Schema.org)      & 6 detail page templates \\
Microformats 2            & Structured data in HTML (h-card, h-event) & 6 detail page templates \\
mf2py                     & Server-side MF2 parsing                    & \texttt{/api/parse-microformats/} \\
\bottomrule
\end{tabular}
\end{center}

\subsection{Dataset}

We used the \textbf{Kaggle Formula 1 World Championship} dataset (1950--2024),
which follows the Ergast-style CSV schema.
It covers 75 seasons, 861 drivers, 212 constructors, 77 circuits, and approximately 25{,}000 race results
distributed across 13 CSV files. These are loaded into a GraphDB repository named \texttt{ws-formula1-owlmax}.

% =============================================================================
\section{Ontology Definition (RDFS and OWL)}
% =============================================================================

We designed the ontology as a single Turtle file (\texttt{data/rdf/formula1\_ontology.ttl}),
following the OWL 2 DL profile. It is the \textbf{only hand-authored source file}
in the data pipeline --- every other RDF file is either generated by a script or
produced by the reasoner.

\subsection{OWL Profile}

We chose \textbf{OWL 2 DL} because it is a decidable fragment of first-order logic
with well-defined reasoning guarantees, and because GraphDB supports the
\textbf{OWL-Max with RDFS} ruleset out of the box. This ruleset applies
\texttt{rdfs:domain}, \texttt{rdfs:range}, \texttt{rdfs:subClassOf},
\texttt{owl:inverseOf}, and \texttt{owl:SymmetricProperty} entailments at triple
load time --- meaning the reasoner fires before any SPARQL query runs, and the
inferred triples are already in the graph when the application queries it.

\subsection{Class Hierarchy}

We defined 20 asserted OWL classes (plus 7 SPIN-inferred subclasses added at runtime)
organised into a domain-aligned hierarchy. A complete table of all 27 classes with their
sources is provided in Section~\ref{sec:metrics}. The key structural excerpt is below:

\begin{lstlisting}[style=turtle, caption={Top-level class hierarchy in the ontology}]
f1:Entity      rdfs:subClassOf  owl:Thing .
f1:Person      rdfs:subClassOf  f1:Entity .
f1:Team        rdfs:subClassOf  f1:Entity .
f1:Venue       rdfs:subClassOf  f1:Entity .
f1:Driver      rdfs:subClassOf  f1:Person .
f1:Constructor rdfs:subClassOf  f1:Team .
f1:Circuit     rdfs:subClassOf  f1:Venue .
f1:Race        rdfs:subClassOf  f1:Event .
f1:Season      rdfs:subClassOf  f1:Event .
f1:Result      rdfs:subClassOf  f1:PerformanceRecord .
% Inferred subclasses (populated by SPIN rules):
f1:WorldChampion      rdfs:subClassOf  f1:Driver .
f1:MultichampionDriver rdfs:subClassOf f1:WorldChampion .
f1:Veteran            rdfs:subClassOf  f1:Driver .
f1:ConstructorChampion rdfs:subClassOf f1:Constructor .
f1:ActiveCircuit      rdfs:subClassOf  f1:Circuit .
f1:HistoricCircuit    rdfs:subClassOf  f1:Circuit .
f1:PodiumResult       rdfs:subClassOf  f1:Result .
\end{lstlisting}

\subsection{Base Class Inference Without Explicit rdf:type}

A central design requirement was that the facts file should contain \textbf{no explicit}
\texttt{rdf:type f1:Driver} (or Constructor, Circuit) assertions. The ontology must
infer these types purely from the properties present on each entity.

This is achieved using two complementary mechanisms:

\subsubsection{rdfs:domain (for GraphDB)}

Each ``signature'' property that appears exclusively on one entity class is declared
with \texttt{rdfs:domain}. GraphDB's OWL-Max ruleset then infers \texttt{rdf:type} at
load time from any occurrence of that property:

\begin{lstlisting}[style=turtle, caption={rdfs:domain declarations as inference anchors}]
f1:driverRef       rdfs:domain  f1:Driver .
f1:constructorRef  rdfs:domain  f1:Constructor .
f1:circuitRef      rdfs:domain  f1:Circuit .
\end{lstlisting}

So \texttt{<driver/1> f1:driverRef "hamilton"} causes GraphDB to automatically
infer \texttt{<driver/1> rdf:type f1:Driver} without any explicit assertion.

\subsubsection{owl:equivalentClass (for Protégé / HermiT)}

For validation in Protégé, the OWL reasoner HermiT requires a
\emph{necessary and sufficient} condition. Each base class is therefore defined with
\texttt{owl:equivalentClass} using a \texttt{someValuesFrom} restriction:

\begin{lstlisting}[style=turtle, caption={Bidirectional equivalentClass definition for Driver}]
f1:Driver  owl:equivalentClass  [
    rdf:type        owl:Restriction ;
    owl:onProperty  f1:driverRef ;
    owl:someValuesFrom  xsd:string
] .
\end{lstlisting}

This means: an individual is a \texttt{f1:Driver} \textbf{if and only if} it has
at least one \texttt{f1:driverRef} value. HermiT can classify anonymous individuals
without an explicit type assertion, making the ontology verifiable in Protégé.

\begin{infobox}
\textbf{Verification:} Loading the facts file alone into GraphDB (without the
ontology) returns 0 results for \texttt{SELECT * WHERE \{?d rdf:type f1:Driver\}}.
After loading the integrated file (ontology + facts), the same query returns 861 drivers.
The ontology, not the data, does the classification.
\end{infobox}

\subsection{Why rdfs:domain is Omitted on Foreign-Key Properties}

Properties like \texttt{f1:driverId}, \texttt{f1:constructorId}, and \texttt{f1:raceId}
appear as \textbf{literal columns on multiple entity types}. For example, a
\texttt{Result} row carries both a \texttt{f1:driverId} and a \texttt{f1:constructorId}
literal. Declaring \texttt{rdfs:domain f1:Driver} on \texttt{f1:driverId} would cause
GraphDB to infer that every \texttt{Result} individual is also a \texttt{Driver},
immediately triggering the \texttt{owl:disjointWith f1:Constructor} inconsistency.
Domain declarations are intentionally omitted for these FK properties.

\subsection{Key OWL Axioms}

\begin{center}
\begin{tabularx}{\linewidth}{@{}llX@{}}
\toprule
\textbf{Construct} & \textbf{Applied to} & \textbf{Purpose} \\
\midrule
\texttt{owl:FunctionalProperty} & \texttt{race}, \texttt{driver}, \texttt{status},
  \texttt{circuit}, \texttt{season}, \texttt{heldAt}
  & At most one value per subject \\[0.4em]
\texttt{owl:SymmetricProperty} & \texttt{wasTeammate}
  & $x$ wasTeammate $y \Rightarrow y$ wasTeammate $x$ (GraphDB materialises the inverse) \\[0.4em]
\texttt{owl:inverseOf} & \texttt{hadDriver} $\leftrightarrow$ \texttt{drovFor}
  & Constructor $\to$ Driver navigation without storing the triple twice \\[0.4em]
\texttt{rdfs:subPropertyOf} & \texttt{wonRace} $\sqsubseteq$ \texttt{competedIn},
  \texttt{achievedPodium} $\sqsubseteq$ \texttt{competedIn}
  & Race wins also appear as \texttt{competedIn} triples; enables property-path queries \\[0.4em]
\texttt{owl:maxCardinality 1} & \texttt{Race $\sqsubseteq$ $\leq$1 circuit}
  & Each race is held at exactly one circuit \\[0.4em]
\texttt{owl:disjointWith} & Driver/Constructor, Driver/Circuit, Race/Season,
  Result/Qualifying\-Result, \ldots
  & Detects modelling errors during Protégé validation \\
\bottomrule
\end{tabularx}
\end{center}

\noindent Complete tables of all 27 classes and 21 object properties are provided
in Section~\ref{sec:metrics} (\textit{Quantitative Analysis}).

\subsection{Restriction Axioms for Inferred Subclasses}

Three inferred subclasses carry \texttt{rdfs:subClassOf} restriction axioms
that document the structural intent --- even though the actual classification
is performed by SPIN rules (Section~\ref{sec:spin}):

\begin{lstlisting}[style=turtle, caption={Restriction on WorldChampion: must have at least one wonChampionship}]
f1:WorldChampion  rdfs:subClassOf  [
    rdf:type          owl:Restriction ;
    owl:onProperty    f1:wonChampionship ;
    owl:someValuesFrom  f1:Season
] .
\end{lstlisting}

\subsection{Inference Verification}

After loading the integrated file (389 ontology triples + 6.6M fact triples) and
running all 21 SPIN rules, all inference checks pass:

\begin{center}
\begin{tabular}{@{}llr@{}}
\toprule
\textbf{Mechanism} & \textbf{Inferred} & \textbf{Count} \\
\midrule
\texttt{rdfs:domain driverRef} & \texttt{f1:Driver} entities & 861 \\
\texttt{rdfs:domain constructorRef} & \texttt{f1:Constructor} entities & 212 \\
\texttt{rdfs:domain circuitRef} & \texttt{f1:Circuit} entities & 77 \\
SPIN \texttt{classify\_WorldChampion} & \texttt{f1:WorldChampion} & 34 \\
SPIN \texttt{classify\_MultichampionDriver} & \texttt{f1:MultichampionDriver} & 17 \\
SPIN \texttt{classify\_Veteran} & \texttt{f1:Veteran} & 86 \\
SPIN \texttt{classify\_ConstructorChampion} & \texttt{f1:ConstructorChampion} & 17 \\
SPIN \texttt{classify\_ActiveCircuit} & \texttt{f1:ActiveCircuit} & 24 \\
SPIN \texttt{classify\_HistoricCircuit} & \texttt{f1:HistoricCircuit} & 45 \\
SPIN \texttt{classify\_PodiumResult} & \texttt{f1:PodiumResult} & 3{,}397 \\
SPIN \texttt{infer\_wonRace} & \texttt{f1:wonRace} triples & 1{,}128 \\
SPIN \texttt{infer\_startedFromP1} & \texttt{f1:startedFromP1} triples & 1{,}135 \\
SPIN \texttt{infer\_convertedP1ToWin} & \texttt{f1:convertedP1ToWin} triples & 486 \\
SPIN \texttt{infer\_setFastestLap} & \texttt{f1:setFastestLap} triples & 410 \\
SPIN \texttt{infer\_achievedHatTrick} & \texttt{f1:achievedHatTrick} triples & 69 \\
SPIN \texttt{infer\_wasTeammate} & \texttt{f1:wasTeammate} triples & 10{,}012 \\
SPIN \texttt{infer\_wonChampionship} & \texttt{f1:wonChampionship} triples & 75 \\
\texttt{owl:inverseOf} & \texttt{f1:hadDriver} triples & 2{,}150 \\
SPIN \texttt{infer\_heldAt} & \texttt{f1:heldAt} triples & 1{,}125 \\
\bottomrule
\end{tabular}
\end{center}

% =============================================================================
\subsection{Protégé Validation with HermiT Reasoner}
\label{sec:hermit}
% =============================================================================

The ontology was formally validated using the \textbf{HermiT 1.4.3} DL reasoner
(bundled with Protégé 5.6.9) running programmatically via the \texttt{owlready2}
Python library. This section documents the complete analysis, including a real
issue encountered and the design decision taken to resolve it.

\subsubsection{OWL 2 Datatype Map Compliance}

During validation, HermiT raised an \texttt{UnsupportedDatatypeException} for three
\texttt{rdfs:range} declarations:

\begin{lstlisting}[style=turtle, caption={Original range declarations causing HermiT rejection}]
f1:dob   rdfs:range  xsd:date .   -- NOT in OWL 2 datatype map
f1:date  rdfs:range  xsd:date .   -- NOT in OWL 2 datatype map
f1:year  rdfs:range  xsd:gYear .  -- NOT in OWL 2 datatype map
\end{lstlisting}

The OWL 2 specification defines a strict \emph{datatype map} --- a subset of XML
Schema datatypes that OWL 2 DL reasoners are required to support. This map includes
\texttt{xsd:dateTime} and \texttt{xsd:dateTimeStamp} but \textbf{not}
\texttt{xsd:date} or \texttt{xsd:gYear}. The teacher's slides state that
``A OWL faz uso dos tipos de dados definidos no XML Schema'' (OWL uses XML Schema
datatypes), but the key constraint is that only a \emph{subset} of XML Schema
types are in the OWL 2 DL datatype map.

\subsubsection{Design Decision}

The \texttt{rdfs:range} declaration is defined in RDFS as \emph{optional}: it
states what kind of values a property \emph{may} take, but its absence does not
prevent the property from functioning. The decision was to
\textbf{remove the three non-compliant range declarations} rather than replacing
them with an incorrect type (e.g., \texttt{xsd:dateTime}) or a loss of type
information (e.g., \texttt{xsd:string}).

\textbf{Justification:} The semantic power of the ontology is entirely unchanged.
All inference mechanisms --- \texttt{rdfs:domain} for base class classification,
\texttt{owl:equivalentClass} for Driver/Constructor/Circuit,
\texttt{owl:disjointWith} for consistency checking,
\texttt{owl:FunctionalProperty}, \texttt{owl:SymmetricProperty},
\texttt{owl:inverseOf}, and \texttt{rdfs:subPropertyOf} --- are independent of
these three range declarations. The domain declarations
(\texttt{f1:dob rdfs:domain f1:Driver}, \texttt{f1:year rdfs:domain f1:Event},
\texttt{f1:date rdfs:domain f1:Race}) are preserved and still provide useful
type inference. The actual data values in GraphDB remain typed as
\texttt{xsd:date} and \texttt{xsd:gYear}, which GraphDB handles correctly
under its OWL-Max ruleset.

\subsubsection{HermiT Verification Results}

HermiT was run twice: once on the original ontology (838~ms), and again after
adding the two new properties (\texttt{f1:finishedSecond}, \texttt{f1:finishedThird})
and restructuring the subPropertyOf chain (\texttt{wonRace}~$\sqsubseteq$~\texttt{achievedPodium}).
Both runs produced identical consistency results. The updated ontology (389 triples,
27 classes, 21 object properties) completed reasoning in \textbf{1{,}333~ms}.

\begin{center}
\begin{tabular}{@{}lp{10cm}@{}}
\toprule
\textbf{Check} & \textbf{Result} \\
\midrule
Consistency & \textbf{CONSISTENT} --- no exception, no contradiction \\
Unsatisfiable classes & \textbf{0} --- every named class has a non-empty extension \\
Equivalent class definitions & Driver~$\equiv$~$\exists$\texttt{driverRef.string},
  Constructor~$\equiv$~$\exists$\texttt{constructorRef.string},
  Circuit~$\equiv$~$\exists$\texttt{circuitRef.string} \\
Property hierarchy & \texttt{wonRace}~$\sqsubseteq$~\texttt{achievedPodium}~$\sqsubseteq$~\texttt{competedIn},
  \texttt{finishedSecond}~$\sqsubseteq$~\texttt{achievedPodium},
  \texttt{finishedThird}~$\sqsubseteq$~\texttt{achievedPodium} \\
Functional properties & \texttt{circuit}, \texttt{driver}, \texttt{heldAt},
  \texttt{race}, \texttt{season}, \texttt{status} \\
Symmetric properties & \texttt{wasTeammate} \\
Disjoint classes & 7 pairs verified (Driver~$\perp$~Constructor, Driver~$\perp$~Circuit,
  Circuit~$\perp$~Constructor, Race~$\perp$~Season, Driver~$\perp$~Race,
  Result~$\perp$~QualifyingResult, LapTime~$\perp$~Result) \\
\bottomrule
\end{tabular}
\end{center}

The full inferred class hierarchy produced by HermiT:

\begin{lstlisting}[style=shell, caption={Class hierarchy inferred by HermiT from owl:equivalentClass + rdfs:subClassOf axioms}]
owl:Thing
  Entity
    Person  ->  Driver  ->  WorldChampion  ->  MultichampionDriver
                         ->  Veteran
    Team    ->  Constructor  ->  ConstructorChampion
    Venue   ->  Circuit      ->  ActiveCircuit
                             ->  HistoricCircuit
  Event
    Race
    Season
  PerformanceRecord
    Result  ->  PodiumResult
    QualifyingResult
    LapTime
    PitStop
    ConstructorResult
  Standing
    DriverStanding
    ConstructorStanding
  Status
\end{lstlisting}

This hierarchy proves the ontology's inference power: the
\texttt{WorldChampion}, \texttt{Veteran}, \texttt{ConstructorChampion},
\texttt{ActiveCircuit}, \texttt{HistoricCircuit}, and \texttt{PodiumResult}
subclasses are all correctly positioned by HermiT using
\texttt{rdfs:subClassOf} alone, confirming that the ontology structure is
semantically sound.

% =============================================================================
\section{Quantitative Analysis}
\label{sec:metrics}
% =============================================================================

This section provides objective, verifiable evidence of the semantic contribution
made by the ontology and inference system.

\subsection{Ontology Metrics}

% LaTeX table: Quantitative metrics
% Include in report with: % LaTeX table: Quantitative metrics
% Include in report with: % LaTeX table: Quantitative metrics
% Include in report with: \input{tables/metrics.tex}

% ── Ontology metrics ──────────────────────────────────────────────────────
\begin{center}
\begin{tabular}{@{}lrl@{}}
\toprule
\textbf{Ontology metric} & \textbf{Count} & \textbf{Notes} \\
\midrule
OWL classes (total)          & 27  & 7 abstract + 13 concrete + 7 SPIN-inferred \\
Concrete asserted classes    & 13  & Populated from CSV \\
SPIN-inferred subclasses     & 7   & Not in facts file \\
Object properties            & 21  & 6 raw + 15 SPIN-materialised \\
Datatype properties          & 35  & Covering all CSV columns \\
owl:FunctionalProperty       & 6   & race, driver, circuit, status, season, heldAt \\
owl:SymmetricProperty        & 1   & f1:wasTeammate \\
owl:inverseOf pairs          & 1   & hadDriver / drovFor \\
rdfs:subPropertyOf chains    & 6   & All under f1:competedIn \\
owl:disjointWith pairs       & 7   & Prevents cross-type inference \\
owl:maxCardinality restrictions & 1 & Race$\to$circuit (max 1) \\
owl:equivalentClass (Protégé)& 3   & Driver, Constructor, Circuit base inference \\
owl:someValuesFrom restrictions & 2 & WorldChampion, ConstructorChampion \\
HermiT reasoning time        & 1{,}333 ms & owlready2 on Python 3.12 (21 properties) \\
Unsatisfiable classes (HermiT) & 0 & Ontology is consistent \\
\bottomrule
\end{tabular}
\end{center}

\vspace{0.8em}

% ── GraphDB / data pipeline metrics ──────────────────────────────────────
\begin{center}
\begin{tabular}{@{}lrl@{}}
\toprule
\textbf{Data pipeline metric} & \textbf{Count} & \textbf{Notes} \\
\midrule
Raw CSV files                & 13  & Ergast F1 1950--2024 \\
N-Triple facts               & 6{,}627{,}045 & formula1.nt \\
Ontology triples             & 375 & formula1\_ontology.ttl \\
Integrated file triples      & 6{,}627{,}420 & formula1\_integrated.nt \\
OWL-Max inferred (at import) & $\approx$ 1{,}950 & rdfs:domain $\to$ rdf:type \\
SPIN-generated triples       & 22{,}921 & 21 rules (verified live, June 2026) \\
Total triples in repository  & 6{,}650{,}355 & post-inference (verified live) \\
\midrule
Drivers                      & 861 & OWL-Max inferred \\
Constructors                 & 212 & OWL-Max inferred \\
Circuits                     & 77  & OWL-Max inferred \\
Seasons                      & 75  & Explicit rdf:type \\
Races                        & 1{,}125 & Explicit rdf:type \\
Results                      & $\approx$ 25{,}000 & Explicit rdf:type \\
\bottomrule
\end{tabular}
\end{center}

\vspace{0.8em}

% ── Semantic markup coverage ──────────────────────────────────────────────
\begin{center}
\begin{tabularx}{\linewidth}{@{}lllX@{}}
\toprule
\textbf{Page} & \textbf{RDFa typeof} & \textbf{MF2 classes} & \textbf{Properties} \\
\midrule
Driver detail      & schema:Person                 & h-card        & name, nationality, birthDate, sameAs \\
Constructor detail & schema:SportsOrganization     & h-card        & name, nationality, sameAs \\
Circuit detail     & schema:SportsActivityLocation & h-card, h-adr & name, addressLocality, addressCountry, lat, lng \\
Race detail        & schema:SportsEvent            & h-event       & name, location, startDate, sameAs \\
Season detail      & schema:SportsEvent            & h-event       & name, startDate, sameAs \\
Champions list     & schema:ItemList               & ---           & itemListElement (Person $\times$ 34) \\
Homepage           & schema:Dataset                & ---           & name, description, license, temporalCoverage \\
\bottomrule
\end{tabularx}
\end{center}


% ── Ontology metrics ──────────────────────────────────────────────────────
\begin{center}
\begin{tabular}{@{}lrl@{}}
\toprule
\textbf{Ontology metric} & \textbf{Count} & \textbf{Notes} \\
\midrule
OWL classes (total)          & 27  & 7 abstract + 13 concrete + 7 SPIN-inferred \\
Concrete asserted classes    & 13  & Populated from CSV \\
SPIN-inferred subclasses     & 7   & Not in facts file \\
Object properties            & 21  & 6 raw + 15 SPIN-materialised \\
Datatype properties          & 35  & Covering all CSV columns \\
owl:FunctionalProperty       & 6   & race, driver, circuit, status, season, heldAt \\
owl:SymmetricProperty        & 1   & f1:wasTeammate \\
owl:inverseOf pairs          & 1   & hadDriver / drovFor \\
rdfs:subPropertyOf chains    & 6   & All under f1:competedIn \\
owl:disjointWith pairs       & 7   & Prevents cross-type inference \\
owl:maxCardinality restrictions & 1 & Race$\to$circuit (max 1) \\
owl:equivalentClass (Protégé)& 3   & Driver, Constructor, Circuit base inference \\
owl:someValuesFrom restrictions & 2 & WorldChampion, ConstructorChampion \\
HermiT reasoning time        & 1{,}333 ms & owlready2 on Python 3.12 (21 properties) \\
Unsatisfiable classes (HermiT) & 0 & Ontology is consistent \\
\bottomrule
\end{tabular}
\end{center}

\vspace{0.8em}

% ── GraphDB / data pipeline metrics ──────────────────────────────────────
\begin{center}
\begin{tabular}{@{}lrl@{}}
\toprule
\textbf{Data pipeline metric} & \textbf{Count} & \textbf{Notes} \\
\midrule
Raw CSV files                & 13  & Ergast F1 1950--2024 \\
N-Triple facts               & 6{,}627{,}045 & formula1.nt \\
Ontology triples             & 375 & formula1\_ontology.ttl \\
Integrated file triples      & 6{,}627{,}420 & formula1\_integrated.nt \\
OWL-Max inferred (at import) & $\approx$ 1{,}950 & rdfs:domain $\to$ rdf:type \\
SPIN-generated triples       & 22{,}921 & 21 rules (verified live, June 2026) \\
Total triples in repository  & 6{,}650{,}355 & post-inference (verified live) \\
\midrule
Drivers                      & 861 & OWL-Max inferred \\
Constructors                 & 212 & OWL-Max inferred \\
Circuits                     & 77  & OWL-Max inferred \\
Seasons                      & 75  & Explicit rdf:type \\
Races                        & 1{,}125 & Explicit rdf:type \\
Results                      & $\approx$ 25{,}000 & Explicit rdf:type \\
\bottomrule
\end{tabular}
\end{center}

\vspace{0.8em}

% ── Semantic markup coverage ──────────────────────────────────────────────
\begin{center}
\begin{tabularx}{\linewidth}{@{}lllX@{}}
\toprule
\textbf{Page} & \textbf{RDFa typeof} & \textbf{MF2 classes} & \textbf{Properties} \\
\midrule
Driver detail      & schema:Person                 & h-card        & name, nationality, birthDate, sameAs \\
Constructor detail & schema:SportsOrganization     & h-card        & name, nationality, sameAs \\
Circuit detail     & schema:SportsActivityLocation & h-card, h-adr & name, addressLocality, addressCountry, lat, lng \\
Race detail        & schema:SportsEvent            & h-event       & name, location, startDate, sameAs \\
Season detail      & schema:SportsEvent            & h-event       & name, startDate, sameAs \\
Champions list     & schema:ItemList               & ---           & itemListElement (Person $\times$ 34) \\
Homepage           & schema:Dataset                & ---           & name, description, license, temporalCoverage \\
\bottomrule
\end{tabularx}
\end{center}


% ── Ontology metrics ──────────────────────────────────────────────────────
\begin{center}
\begin{tabular}{@{}lrl@{}}
\toprule
\textbf{Ontology metric} & \textbf{Count} & \textbf{Notes} \\
\midrule
OWL classes (total)          & 27  & 7 abstract + 13 concrete + 7 SPIN-inferred \\
Concrete asserted classes    & 13  & Populated from CSV \\
SPIN-inferred subclasses     & 7   & Not in facts file \\
Object properties            & 21  & 6 raw + 15 SPIN-materialised \\
Datatype properties          & 35  & Covering all CSV columns \\
owl:FunctionalProperty       & 6   & race, driver, circuit, status, season, heldAt \\
owl:SymmetricProperty        & 1   & f1:wasTeammate \\
owl:inverseOf pairs          & 1   & hadDriver / drovFor \\
rdfs:subPropertyOf chains    & 6   & All under f1:competedIn \\
owl:disjointWith pairs       & 7   & Prevents cross-type inference \\
owl:maxCardinality restrictions & 1 & Race$\to$circuit (max 1) \\
owl:equivalentClass (Protégé)& 3   & Driver, Constructor, Circuit base inference \\
owl:someValuesFrom restrictions & 2 & WorldChampion, ConstructorChampion \\
HermiT reasoning time        & 1{,}333 ms & owlready2 on Python 3.12 (21 properties) \\
Unsatisfiable classes (HermiT) & 0 & Ontology is consistent \\
\bottomrule
\end{tabular}
\end{center}

\vspace{0.8em}

% ── GraphDB / data pipeline metrics ──────────────────────────────────────
\begin{center}
\begin{tabular}{@{}lrl@{}}
\toprule
\textbf{Data pipeline metric} & \textbf{Count} & \textbf{Notes} \\
\midrule
Raw CSV files                & 13  & Ergast F1 1950--2024 \\
N-Triple facts               & 6{,}627{,}045 & formula1.nt \\
Ontology triples             & 375 & formula1\_ontology.ttl \\
Integrated file triples      & 6{,}627{,}420 & formula1\_integrated.nt \\
OWL-Max inferred (at import) & $\approx$ 1{,}950 & rdfs:domain $\to$ rdf:type \\
SPIN-generated triples       & 22{,}921 & 21 rules (verified live, June 2026) \\
Total triples in repository  & 6{,}650{,}355 & post-inference (verified live) \\
\midrule
Drivers                      & 861 & OWL-Max inferred \\
Constructors                 & 212 & OWL-Max inferred \\
Circuits                     & 77  & OWL-Max inferred \\
Seasons                      & 75  & Explicit rdf:type \\
Races                        & 1{,}125 & Explicit rdf:type \\
Results                      & $\approx$ 25{,}000 & Explicit rdf:type \\
\bottomrule
\end{tabular}
\end{center}

\vspace{0.8em}

% ── Semantic markup coverage ──────────────────────────────────────────────
\begin{center}
\begin{tabularx}{\linewidth}{@{}lllX@{}}
\toprule
\textbf{Page} & \textbf{RDFa typeof} & \textbf{MF2 classes} & \textbf{Properties} \\
\midrule
Driver detail      & schema:Person                 & h-card        & name, nationality, birthDate, sameAs \\
Constructor detail & schema:SportsOrganization     & h-card        & name, nationality, sameAs \\
Circuit detail     & schema:SportsActivityLocation & h-card, h-adr & name, addressLocality, addressCountry, lat, lng \\
Race detail        & schema:SportsEvent            & h-event       & name, location, startDate, sameAs \\
Season detail      & schema:SportsEvent            & h-event       & name, startDate, sameAs \\
Champions list     & schema:ItemList               & ---           & itemListElement (Person $\times$ 34) \\
Homepage           & schema:Dataset                & ---           & name, description, license, temporalCoverage \\
\bottomrule
\end{tabularx}
\end{center}


\subsection{Complete Class Inventory}

% LaTeX table: Ontology classes
% Include in report with: % LaTeX table: Ontology classes
% Include in report with: % LaTeX table: Ontology classes
% Include in report with: \input{tables/classes.tex}
\begin{center}
\begin{tabular}{@{}llll@{}}
\toprule
\textbf{Class} & \textbf{Superclass} & \textbf{Type} & \textbf{Source} \\
\midrule
\texttt{f1:Entity}              & —                      & Abstract root     & Hand-authored \\
\texttt{f1:Person}              & \texttt{f1:Entity}     & Abstract          & Hand-authored \\
\texttt{f1:Team}                & \texttt{f1:Entity}     & Abstract          & Hand-authored \\
\texttt{f1:Venue}               & \texttt{f1:Entity}     & Abstract          & Hand-authored \\
\texttt{f1:Event}               & —                      & Abstract          & Hand-authored \\
\texttt{f1:PerformanceRecord}   & —                      & Abstract          & Hand-authored \\
\texttt{f1:Standing}            & —                      & Abstract          & Hand-authored \\
\midrule
\texttt{f1:Driver}              & \texttt{f1:Person}     & Concrete; asserted  & OWL-Max (rdfs:domain) \\
\texttt{f1:Constructor}         & \texttt{f1:Team}       & Concrete; asserted  & OWL-Max (rdfs:domain) \\
\texttt{f1:Circuit}             & \texttt{f1:Venue}      & Concrete; asserted  & OWL-Max (rdfs:domain) \\
\texttt{f1:Season}              & \texttt{f1:Event}      & Concrete; asserted  & Explicit in facts \\
\texttt{f1:Race}                & \texttt{f1:Event}      & Concrete; asserted  & Explicit in facts \\
\texttt{f1:Result}              & \texttt{f1:PerformanceRecord} & Concrete   & Explicit in facts \\
\texttt{f1:QualifyingResult}    & \texttt{f1:PerformanceRecord} & Concrete   & Explicit in facts \\
\texttt{f1:LapTime}             & \texttt{f1:PerformanceRecord} & Concrete   & Explicit in facts \\
\texttt{f1:PitStop}             & \texttt{f1:PerformanceRecord} & Concrete   & Explicit in facts \\
\texttt{f1:ConstructorResult}   & \texttt{f1:PerformanceRecord} & Concrete   & Explicit in facts \\
\texttt{f1:DriverStanding}      & \texttt{f1:Standing}   & Concrete; asserted  & Explicit in facts \\
\texttt{f1:ConstructorStanding} & \texttt{f1:Standing}   & Concrete; asserted  & Explicit in facts \\
\texttt{f1:Status}              & —                      & Concrete; asserted  & Explicit in facts \\
\midrule
\texttt{f1:WorldChampion}       & \texttt{f1:Driver}     & \textbf{SPIN-inferred} & 34 instances \\
\texttt{f1:MultichampionDriver} & \texttt{f1:WorldChampion} & \textbf{SPIN-inferred} & 17 instances \\
\texttt{f1:Veteran}             & \texttt{f1:Driver}     & \textbf{SPIN-inferred} & 86 instances \\
\texttt{f1:ActiveCircuit}       & \texttt{f1:Circuit}    & \textbf{SPIN-inferred} & 24 instances \\
\texttt{f1:HistoricCircuit}     & \texttt{f1:Circuit}    & \textbf{SPIN-inferred} & 45 instances \\
\texttt{f1:ConstructorChampion} & \texttt{f1:Constructor}& \textbf{SPIN-inferred} & 17 instances \\
\texttt{f1:PodiumResult}        & \texttt{f1:Result}     & \textbf{SPIN-inferred} & 3{,}397 instances \\
\bottomrule
\end{tabular}
\end{center}

\begin{center}
\begin{tabular}{@{}llll@{}}
\toprule
\textbf{Class} & \textbf{Superclass} & \textbf{Type} & \textbf{Source} \\
\midrule
\texttt{f1:Entity}              & —                      & Abstract root     & Hand-authored \\
\texttt{f1:Person}              & \texttt{f1:Entity}     & Abstract          & Hand-authored \\
\texttt{f1:Team}                & \texttt{f1:Entity}     & Abstract          & Hand-authored \\
\texttt{f1:Venue}               & \texttt{f1:Entity}     & Abstract          & Hand-authored \\
\texttt{f1:Event}               & —                      & Abstract          & Hand-authored \\
\texttt{f1:PerformanceRecord}   & —                      & Abstract          & Hand-authored \\
\texttt{f1:Standing}            & —                      & Abstract          & Hand-authored \\
\midrule
\texttt{f1:Driver}              & \texttt{f1:Person}     & Concrete; asserted  & OWL-Max (rdfs:domain) \\
\texttt{f1:Constructor}         & \texttt{f1:Team}       & Concrete; asserted  & OWL-Max (rdfs:domain) \\
\texttt{f1:Circuit}             & \texttt{f1:Venue}      & Concrete; asserted  & OWL-Max (rdfs:domain) \\
\texttt{f1:Season}              & \texttt{f1:Event}      & Concrete; asserted  & Explicit in facts \\
\texttt{f1:Race}                & \texttt{f1:Event}      & Concrete; asserted  & Explicit in facts \\
\texttt{f1:Result}              & \texttt{f1:PerformanceRecord} & Concrete   & Explicit in facts \\
\texttt{f1:QualifyingResult}    & \texttt{f1:PerformanceRecord} & Concrete   & Explicit in facts \\
\texttt{f1:LapTime}             & \texttt{f1:PerformanceRecord} & Concrete   & Explicit in facts \\
\texttt{f1:PitStop}             & \texttt{f1:PerformanceRecord} & Concrete   & Explicit in facts \\
\texttt{f1:ConstructorResult}   & \texttt{f1:PerformanceRecord} & Concrete   & Explicit in facts \\
\texttt{f1:DriverStanding}      & \texttt{f1:Standing}   & Concrete; asserted  & Explicit in facts \\
\texttt{f1:ConstructorStanding} & \texttt{f1:Standing}   & Concrete; asserted  & Explicit in facts \\
\texttt{f1:Status}              & —                      & Concrete; asserted  & Explicit in facts \\
\midrule
\texttt{f1:WorldChampion}       & \texttt{f1:Driver}     & \textbf{SPIN-inferred} & 34 instances \\
\texttt{f1:MultichampionDriver} & \texttt{f1:WorldChampion} & \textbf{SPIN-inferred} & 17 instances \\
\texttt{f1:Veteran}             & \texttt{f1:Driver}     & \textbf{SPIN-inferred} & 86 instances \\
\texttt{f1:ActiveCircuit}       & \texttt{f1:Circuit}    & \textbf{SPIN-inferred} & 24 instances \\
\texttt{f1:HistoricCircuit}     & \texttt{f1:Circuit}    & \textbf{SPIN-inferred} & 45 instances \\
\texttt{f1:ConstructorChampion} & \texttt{f1:Constructor}& \textbf{SPIN-inferred} & 17 instances \\
\texttt{f1:PodiumResult}        & \texttt{f1:Result}     & \textbf{SPIN-inferred} & 3{,}397 instances \\
\bottomrule
\end{tabular}
\end{center}

\begin{center}
\begin{tabular}{@{}llll@{}}
\toprule
\textbf{Class} & \textbf{Superclass} & \textbf{Type} & \textbf{Source} \\
\midrule
\texttt{f1:Entity}              & —                      & Abstract root     & Hand-authored \\
\texttt{f1:Person}              & \texttt{f1:Entity}     & Abstract          & Hand-authored \\
\texttt{f1:Team}                & \texttt{f1:Entity}     & Abstract          & Hand-authored \\
\texttt{f1:Venue}               & \texttt{f1:Entity}     & Abstract          & Hand-authored \\
\texttt{f1:Event}               & —                      & Abstract          & Hand-authored \\
\texttt{f1:PerformanceRecord}   & —                      & Abstract          & Hand-authored \\
\texttt{f1:Standing}            & —                      & Abstract          & Hand-authored \\
\midrule
\texttt{f1:Driver}              & \texttt{f1:Person}     & Concrete; asserted  & OWL-Max (rdfs:domain) \\
\texttt{f1:Constructor}         & \texttt{f1:Team}       & Concrete; asserted  & OWL-Max (rdfs:domain) \\
\texttt{f1:Circuit}             & \texttt{f1:Venue}      & Concrete; asserted  & OWL-Max (rdfs:domain) \\
\texttt{f1:Season}              & \texttt{f1:Event}      & Concrete; asserted  & Explicit in facts \\
\texttt{f1:Race}                & \texttt{f1:Event}      & Concrete; asserted  & Explicit in facts \\
\texttt{f1:Result}              & \texttt{f1:PerformanceRecord} & Concrete   & Explicit in facts \\
\texttt{f1:QualifyingResult}    & \texttt{f1:PerformanceRecord} & Concrete   & Explicit in facts \\
\texttt{f1:LapTime}             & \texttt{f1:PerformanceRecord} & Concrete   & Explicit in facts \\
\texttt{f1:PitStop}             & \texttt{f1:PerformanceRecord} & Concrete   & Explicit in facts \\
\texttt{f1:ConstructorResult}   & \texttt{f1:PerformanceRecord} & Concrete   & Explicit in facts \\
\texttt{f1:DriverStanding}      & \texttt{f1:Standing}   & Concrete; asserted  & Explicit in facts \\
\texttt{f1:ConstructorStanding} & \texttt{f1:Standing}   & Concrete; asserted  & Explicit in facts \\
\texttt{f1:Status}              & —                      & Concrete; asserted  & Explicit in facts \\
\midrule
\texttt{f1:WorldChampion}       & \texttt{f1:Driver}     & \textbf{SPIN-inferred} & 34 instances \\
\texttt{f1:MultichampionDriver} & \texttt{f1:WorldChampion} & \textbf{SPIN-inferred} & 17 instances \\
\texttt{f1:Veteran}             & \texttt{f1:Driver}     & \textbf{SPIN-inferred} & 86 instances \\
\texttt{f1:ActiveCircuit}       & \texttt{f1:Circuit}    & \textbf{SPIN-inferred} & 24 instances \\
\texttt{f1:HistoricCircuit}     & \texttt{f1:Circuit}    & \textbf{SPIN-inferred} & 45 instances \\
\texttt{f1:ConstructorChampion} & \texttt{f1:Constructor}& \textbf{SPIN-inferred} & 17 instances \\
\texttt{f1:PodiumResult}        & \texttt{f1:Result}     & \textbf{SPIN-inferred} & 3{,}397 instances \\
\bottomrule
\end{tabular}
\end{center}


\subsection{Complete Object Property Inventory}

% LaTeX table: Object properties
% Include in report with: % LaTeX table: Object properties
% Include in report with: % LaTeX table: Object properties
% Include in report with: \input{tables/object-properties.tex}
\begin{center}
\begin{tabularx}{\linewidth}{@{}lllp{2.2cm}X@{}}
\toprule
\textbf{Property} & \textbf{Domain} & \textbf{Range} & \textbf{OWL axiom} & \textbf{Source} \\
\midrule
\texttt{f1:race}       & (f1:Result) & \texttt{f1:Race}        & Functional & Asserted \\
\texttt{f1:driver}     & (f1:Result) & \texttt{f1:Driver}      & Functional & Asserted \\
\texttt{f1:constructor}& (f1:Result) & \texttt{f1:Constructor} & ---        & Asserted \\
\texttt{f1:status}     & f1:Result   & \texttt{f1:Status}      & Functional & Asserted \\
\texttt{f1:circuit}    & f1:Race     & \texttt{f1:Circuit}     & Functional & Asserted \\
\texttt{f1:season}     & f1:Race     & \texttt{f1:Season}      & Functional & Asserted \\
\midrule
\texttt{f1:heldAt}          & f1:Race        & f1:Circuit     & Functional        & SPIN: infer\_heldAt \\
\texttt{f1:drovFor}         & f1:Driver      & f1:Constructor & ---               & SPIN: infer\_drovFor \\
\texttt{f1:hadDriver}       & f1:Constructor & f1:Driver      & inverseOf         & OWL: inverseOf drovFor \\
\texttt{f1:wasTeammate}     & f1:Driver      & f1:Driver      & Symmetric         & SPIN: infer\_wasTeammate \\
\texttt{f1:competedIn}      & f1:Driver      & f1:Race        & ---               & SPIN (base) \\
\texttt{f1:achievedPodium}  & f1:Driver      & f1:Race        & $\sqsubseteq$ competedIn & SPIN \\
\texttt{f1:finishedSecond}  & f1:Driver      & f1:Race        & $\sqsubseteq$ achievedPodium & SPIN \\
\texttt{f1:finishedThird}   & f1:Driver      & f1:Race        & $\sqsubseteq$ achievedPodium & SPIN \\
\texttt{f1:wonRace}         & f1:Driver      & f1:Race        & $\sqsubseteq$ achievedPodium & SPIN (1{,}128) \\
\texttt{f1:startedFromP1}   & f1:Driver      & f1:Race        & $\sqsubseteq$ competedIn & SPIN (1{,}135) \\
\texttt{f1:setFastestLap}   & f1:Driver      & f1:Race        & $\sqsubseteq$ competedIn & SPIN (410) \\
\texttt{f1:convertedP1ToWin}& f1:Driver      & f1:Race        & $\sqsubseteq$ wonRace     & SPIN (486) \\
\texttt{f1:achievedHatTrick}& f1:Driver      & f1:Race        & $\sqsubseteq$ wonRace     & SPIN (69) \\
\texttt{f1:wonChampionship} & f1:Driver      & f1:Season      & ---               & SPIN (75) \\
\texttt{f1:wonCtorChamp.}   & f1:Constructor & f1:Season      & ---               & SPIN (17) \\
\texttt{f1:finalPoints}     & rdf:Statement  & xsd:decimal    & ---               & Reification metadata \\
\bottomrule
\end{tabularx}
\end{center}

\begin{center}
\begin{tabularx}{\linewidth}{@{}lllp{2.2cm}X@{}}
\toprule
\textbf{Property} & \textbf{Domain} & \textbf{Range} & \textbf{OWL axiom} & \textbf{Source} \\
\midrule
\texttt{f1:race}       & (f1:Result) & \texttt{f1:Race}        & Functional & Asserted \\
\texttt{f1:driver}     & (f1:Result) & \texttt{f1:Driver}      & Functional & Asserted \\
\texttt{f1:constructor}& (f1:Result) & \texttt{f1:Constructor} & ---        & Asserted \\
\texttt{f1:status}     & f1:Result   & \texttt{f1:Status}      & Functional & Asserted \\
\texttt{f1:circuit}    & f1:Race     & \texttt{f1:Circuit}     & Functional & Asserted \\
\texttt{f1:season}     & f1:Race     & \texttt{f1:Season}      & Functional & Asserted \\
\midrule
\texttt{f1:heldAt}          & f1:Race        & f1:Circuit     & Functional        & SPIN: infer\_heldAt \\
\texttt{f1:drovFor}         & f1:Driver      & f1:Constructor & ---               & SPIN: infer\_drovFor \\
\texttt{f1:hadDriver}       & f1:Constructor & f1:Driver      & inverseOf         & OWL: inverseOf drovFor \\
\texttt{f1:wasTeammate}     & f1:Driver      & f1:Driver      & Symmetric         & SPIN: infer\_wasTeammate \\
\texttt{f1:competedIn}      & f1:Driver      & f1:Race        & ---               & SPIN (base) \\
\texttt{f1:achievedPodium}  & f1:Driver      & f1:Race        & $\sqsubseteq$ competedIn & SPIN \\
\texttt{f1:finishedSecond}  & f1:Driver      & f1:Race        & $\sqsubseteq$ achievedPodium & SPIN \\
\texttt{f1:finishedThird}   & f1:Driver      & f1:Race        & $\sqsubseteq$ achievedPodium & SPIN \\
\texttt{f1:wonRace}         & f1:Driver      & f1:Race        & $\sqsubseteq$ achievedPodium & SPIN (1{,}128) \\
\texttt{f1:startedFromP1}   & f1:Driver      & f1:Race        & $\sqsubseteq$ competedIn & SPIN (1{,}135) \\
\texttt{f1:setFastestLap}   & f1:Driver      & f1:Race        & $\sqsubseteq$ competedIn & SPIN (410) \\
\texttt{f1:convertedP1ToWin}& f1:Driver      & f1:Race        & $\sqsubseteq$ wonRace     & SPIN (486) \\
\texttt{f1:achievedHatTrick}& f1:Driver      & f1:Race        & $\sqsubseteq$ wonRace     & SPIN (69) \\
\texttt{f1:wonChampionship} & f1:Driver      & f1:Season      & ---               & SPIN (75) \\
\texttt{f1:wonCtorChamp.}   & f1:Constructor & f1:Season      & ---               & SPIN (17) \\
\texttt{f1:finalPoints}     & rdf:Statement  & xsd:decimal    & ---               & Reification metadata \\
\bottomrule
\end{tabularx}
\end{center}

\begin{center}
\begin{tabularx}{\linewidth}{@{}lllp{2.2cm}X@{}}
\toprule
\textbf{Property} & \textbf{Domain} & \textbf{Range} & \textbf{OWL axiom} & \textbf{Source} \\
\midrule
\texttt{f1:race}       & (f1:Result) & \texttt{f1:Race}        & Functional & Asserted \\
\texttt{f1:driver}     & (f1:Result) & \texttt{f1:Driver}      & Functional & Asserted \\
\texttt{f1:constructor}& (f1:Result) & \texttt{f1:Constructor} & ---        & Asserted \\
\texttt{f1:status}     & f1:Result   & \texttt{f1:Status}      & Functional & Asserted \\
\texttt{f1:circuit}    & f1:Race     & \texttt{f1:Circuit}     & Functional & Asserted \\
\texttt{f1:season}     & f1:Race     & \texttt{f1:Season}      & Functional & Asserted \\
\midrule
\texttt{f1:heldAt}          & f1:Race        & f1:Circuit     & Functional        & SPIN: infer\_heldAt \\
\texttt{f1:drovFor}         & f1:Driver      & f1:Constructor & ---               & SPIN: infer\_drovFor \\
\texttt{f1:hadDriver}       & f1:Constructor & f1:Driver      & inverseOf         & OWL: inverseOf drovFor \\
\texttt{f1:wasTeammate}     & f1:Driver      & f1:Driver      & Symmetric         & SPIN: infer\_wasTeammate \\
\texttt{f1:competedIn}      & f1:Driver      & f1:Race        & ---               & SPIN (base) \\
\texttt{f1:achievedPodium}  & f1:Driver      & f1:Race        & $\sqsubseteq$ competedIn & SPIN \\
\texttt{f1:finishedSecond}  & f1:Driver      & f1:Race        & $\sqsubseteq$ achievedPodium & SPIN \\
\texttt{f1:finishedThird}   & f1:Driver      & f1:Race        & $\sqsubseteq$ achievedPodium & SPIN \\
\texttt{f1:wonRace}         & f1:Driver      & f1:Race        & $\sqsubseteq$ achievedPodium & SPIN (1{,}128) \\
\texttt{f1:startedFromP1}   & f1:Driver      & f1:Race        & $\sqsubseteq$ competedIn & SPIN (1{,}135) \\
\texttt{f1:setFastestLap}   & f1:Driver      & f1:Race        & $\sqsubseteq$ competedIn & SPIN (410) \\
\texttt{f1:convertedP1ToWin}& f1:Driver      & f1:Race        & $\sqsubseteq$ wonRace     & SPIN (486) \\
\texttt{f1:achievedHatTrick}& f1:Driver      & f1:Race        & $\sqsubseteq$ wonRace     & SPIN (69) \\
\texttt{f1:wonChampionship} & f1:Driver      & f1:Season      & ---               & SPIN (75) \\
\texttt{f1:wonCtorChamp.}   & f1:Constructor & f1:Season      & ---               & SPIN (17) \\
\texttt{f1:finalPoints}     & rdf:Statement  & xsd:decimal    & ---               & Reification metadata \\
\bottomrule
\end{tabularx}
\end{center}


\subsection{Ontology Visualizations}
\label{sec:diagrams}

The diagrams below document the ontology and inference architecture. They were
authored in Mermaid and are stored under \texttt{docs/diagrams/}. To regenerate
them as images, run \texttt{mmdc -i docs/diagrams/<file>.md -o report/figures/<name>.png}
(requires the \texttt{@mermaid-js/mermaid-cli} package).

% ── Figure 1: Class hierarchy ────────────────────────────────────────────────
\begin{figure}[H]
\centering
% TODO: replace the fbox below with \includegraphics{figures/ontology-overview}
% after running: mmdc -i docs/diagrams/ontology-overview.md -o report/figures/ontology-overview.png
\fbox{\parbox{0.85\textwidth}{\centering\vspace{1cm}
  \textit{[Figure 1 --- Ontology class hierarchy.\\
   Source: \texttt{docs/diagrams/ontology-overview.md}\\
   Render and replace this placeholder before submission.]}
  \vspace{1cm}}}
\caption{Complete class hierarchy of the F1 ontology, from abstract roots to
SPIN-inferred subclasses. Green nodes are asserted; purple nodes are only
populated by SPIN rules at runtime.}
\label{fig:ontology-hierarchy}
\end{figure}

% ── Figure 2: SPIN inference chain ───────────────────────────────────────────
\begin{figure}[H]
\centering
% TODO: replace with \includegraphics{figures/spin-inference-chain}
% after running: mmdc -i docs/diagrams/spin-inference-chain.md -o report/figures/spin-inference-chain.png
\fbox{\parbox{0.85\textwidth}{\centering\vspace{1cm}
  \textit{[Figure 2 --- SPIN inference chain (3 layers).\\
   Source: \texttt{docs/diagrams/spin-inference-chain.md}\\
   Render and replace this placeholder before submission.]}
  \vspace{1cm}}}
\caption{Three-layer SPIN inference chain: raw CSV data $\to$ Layer~1 race facts
$\to$ Layer~2 combined facts and championships $\to$ Layer~3 entity classifications.}
\label{fig:spin-chain}
\end{figure}

% ── Figure 3: Property hierarchy ─────────────────────────────────────────────
\begin{figure}[H]
\centering
% TODO: replace with \includegraphics{figures/property-hierarchy}
% after running: mmdc -i docs/diagrams/property-hierarchy.md -o report/figures/property-hierarchy.png
\fbox{\parbox{0.85\textwidth}{\centering\vspace{1cm}
  \textit{[Figure 3 --- Object property subPropertyOf hierarchy.\\
   Source: \texttt{docs/diagrams/property-hierarchy.md}\\
   Render and replace this placeholder before submission.]}
  \vspace{1cm}}}
\caption{rdfs:subPropertyOf chain rooted at \texttt{f1:competedIn}. A race win
(\texttt{f1:wonRace}) is also a podium finish (\texttt{f1:achievedPodium}),
which is also a race participation (\texttt{f1:competedIn}).}
\label{fig:property-hierarchy}
\end{figure}

% ── Figure 4: External enrichment ────────────────────────────────────────────
\begin{figure}[H]
\centering
% TODO: replace with \includegraphics{figures/external-enrichment}
% after running: mmdc -i docs/diagrams/external-enrichment.md -o report/figures/external-enrichment.png
\fbox{\parbox{0.85\textwidth}{\centering\vspace{1cm}
  \textit{[Figure 4 --- External knowledge integration architecture.\\
   Source: \texttt{docs/diagrams/external-enrichment.md}\\
   Render and replace this placeholder before submission.]}
  \vspace{1cm}}}
\caption{Architecture of external semantic enrichment: Wikidata (drivers and
constructors), DBpedia (circuits and races), and Wikipedia REST API (images and
extracts). All calls are asynchronous and cached in-process.}
\label{fig:external-enrichment}
\end{figure}

\subsection{Complete SPIN Rule Inventory}

% LaTeX table: All 21 SPIN rules
% Include in report with: % LaTeX table: All 21 SPIN rules
% Include in report with: % LaTeX table: All 21 SPIN rules
% Include in report with: \input{tables/spin-rules.tex}
\begin{center}
\begin{tabularx}{\linewidth}{@{}lXrp{3.5cm}@{}}
\toprule
\textbf{Rule name} & \textbf{Output property / class} & \textbf{Triples} & \textbf{Key pattern} \\
\midrule
\texttt{infer\_wonRace}          & \texttt{f1:wonRace}             & 1{,}128  & positionOrder = 1 \\
\texttt{infer\_finishedSecond}   & \texttt{f1:finishedSecond}      & 1{,}134  & positionOrder = 2 \\
\texttt{infer\_finishedThird}    & \texttt{f1:finishedThird}       & 1{,}135  & positionOrder = 3 \\
\texttt{infer\_achievedPodium}   & \texttt{f1:achievedPodium}      & 3{,}395  & positionOrder $\leq$ 3 \\
\texttt{infer\_startedFromP1}    & \texttt{f1:startedFromP1}       & 1{,}135  & grid = 1 \\
\texttt{infer\_convertedP1ToWin} & \texttt{f1:convertedP1ToWin}   & 486      & startedFromP1 $\wedge$ wonRace \\
\texttt{infer\_setFastestLap}    & \texttt{f1:setFastestLap}       & 410      & rank = 1 \\
\texttt{infer\_achievedHatTrick} & \texttt{f1:achievedHatTrick}    & 69       & wonRace $\wedge$ P1 $\wedge$ fastest \\
\texttt{infer\_drovFor}          & \texttt{f1:drovFor}             & 2{,}150  & result $\to$ driver + constructor \\
\texttt{infer\_wasTeammate}      & \texttt{f1:wasTeammate}         & 10{,}012 & same race + constructor \\
\texttt{infer\_heldAt}           & \texttt{f1:heldAt}              & 1{,}125  & materialise f1:circuit \\
\texttt{infer\_wonChampionship}  & \texttt{f1:wonChampionship}     & 75       & position=1, MAX(round) \\
\texttt{reify\_wonChampionship}  & \texttt{rdf:Statement} nodes    & 75       & annotate with finalPoints \\
\texttt{infer\_wonCtorChamp.}    & \texttt{f1:wonCtorChampionship} & 17       & same pattern \\
\texttt{classify\_WorldChampion} & \texttt{rdf:type f1:WorldChampion}       & 34  & wonChampionship exists \\
\texttt{classify\_Multichamp.}   & \texttt{rdf:type f1:MultichampionDriver} & 17  & COUNT $\geq$ 2 \\
\texttt{classify\_Veteran}       & \texttt{rdf:type f1:Veteran}             & 86  & COUNT(results) $\geq$ 100 \\
\texttt{classify\_ActiveCircuit} & \texttt{rdf:type f1:ActiveCircuit}       & 24  & race in MAX(year) \\
\texttt{classify\_HistoricCircuit}& \texttt{rdf:type f1:HistoricCircuit}    & 45  & NOT EXISTS in last 10 yrs \\
\texttt{classify\_PodiumResult}  & \texttt{rdf:type f1:PodiumResult}        & 3{,}397 & positionOrder $\leq$ 3 \\
\texttt{classify\_CtorChampion}  & \texttt{rdf:type f1:ConstructorChampion} & 17  & wonCtorChampionship \\
\midrule
\multicolumn{2}{l}{\textbf{Total new triples}} & \textbf{22{,}921} & (verified live in GraphDB) \\
\bottomrule
\end{tabularx}
\end{center}

\begin{center}
\begin{tabularx}{\linewidth}{@{}lXrp{3.5cm}@{}}
\toprule
\textbf{Rule name} & \textbf{Output property / class} & \textbf{Triples} & \textbf{Key pattern} \\
\midrule
\texttt{infer\_wonRace}          & \texttt{f1:wonRace}             & 1{,}128  & positionOrder = 1 \\
\texttt{infer\_finishedSecond}   & \texttt{f1:finishedSecond}      & 1{,}134  & positionOrder = 2 \\
\texttt{infer\_finishedThird}    & \texttt{f1:finishedThird}       & 1{,}135  & positionOrder = 3 \\
\texttt{infer\_achievedPodium}   & \texttt{f1:achievedPodium}      & 3{,}395  & positionOrder $\leq$ 3 \\
\texttt{infer\_startedFromP1}    & \texttt{f1:startedFromP1}       & 1{,}135  & grid = 1 \\
\texttt{infer\_convertedP1ToWin} & \texttt{f1:convertedP1ToWin}   & 486      & startedFromP1 $\wedge$ wonRace \\
\texttt{infer\_setFastestLap}    & \texttt{f1:setFastestLap}       & 410      & rank = 1 \\
\texttt{infer\_achievedHatTrick} & \texttt{f1:achievedHatTrick}    & 69       & wonRace $\wedge$ P1 $\wedge$ fastest \\
\texttt{infer\_drovFor}          & \texttt{f1:drovFor}             & 2{,}150  & result $\to$ driver + constructor \\
\texttt{infer\_wasTeammate}      & \texttt{f1:wasTeammate}         & 10{,}012 & same race + constructor \\
\texttt{infer\_heldAt}           & \texttt{f1:heldAt}              & 1{,}125  & materialise f1:circuit \\
\texttt{infer\_wonChampionship}  & \texttt{f1:wonChampionship}     & 75       & position=1, MAX(round) \\
\texttt{reify\_wonChampionship}  & \texttt{rdf:Statement} nodes    & 75       & annotate with finalPoints \\
\texttt{infer\_wonCtorChamp.}    & \texttt{f1:wonCtorChampionship} & 17       & same pattern \\
\texttt{classify\_WorldChampion} & \texttt{rdf:type f1:WorldChampion}       & 34  & wonChampionship exists \\
\texttt{classify\_Multichamp.}   & \texttt{rdf:type f1:MultichampionDriver} & 17  & COUNT $\geq$ 2 \\
\texttt{classify\_Veteran}       & \texttt{rdf:type f1:Veteran}             & 86  & COUNT(results) $\geq$ 100 \\
\texttt{classify\_ActiveCircuit} & \texttt{rdf:type f1:ActiveCircuit}       & 24  & race in MAX(year) \\
\texttt{classify\_HistoricCircuit}& \texttt{rdf:type f1:HistoricCircuit}    & 45  & NOT EXISTS in last 10 yrs \\
\texttt{classify\_PodiumResult}  & \texttt{rdf:type f1:PodiumResult}        & 3{,}397 & positionOrder $\leq$ 3 \\
\texttt{classify\_CtorChampion}  & \texttt{rdf:type f1:ConstructorChampion} & 17  & wonCtorChampionship \\
\midrule
\multicolumn{2}{l}{\textbf{Total new triples}} & \textbf{22{,}921} & (verified live in GraphDB) \\
\bottomrule
\end{tabularx}
\end{center}

\begin{center}
\begin{tabularx}{\linewidth}{@{}lXrp{3.5cm}@{}}
\toprule
\textbf{Rule name} & \textbf{Output property / class} & \textbf{Triples} & \textbf{Key pattern} \\
\midrule
\texttt{infer\_wonRace}          & \texttt{f1:wonRace}             & 1{,}128  & positionOrder = 1 \\
\texttt{infer\_finishedSecond}   & \texttt{f1:finishedSecond}      & 1{,}134  & positionOrder = 2 \\
\texttt{infer\_finishedThird}    & \texttt{f1:finishedThird}       & 1{,}135  & positionOrder = 3 \\
\texttt{infer\_achievedPodium}   & \texttt{f1:achievedPodium}      & 3{,}395  & positionOrder $\leq$ 3 \\
\texttt{infer\_startedFromP1}    & \texttt{f1:startedFromP1}       & 1{,}135  & grid = 1 \\
\texttt{infer\_convertedP1ToWin} & \texttt{f1:convertedP1ToWin}   & 486      & startedFromP1 $\wedge$ wonRace \\
\texttt{infer\_setFastestLap}    & \texttt{f1:setFastestLap}       & 410      & rank = 1 \\
\texttt{infer\_achievedHatTrick} & \texttt{f1:achievedHatTrick}    & 69       & wonRace $\wedge$ P1 $\wedge$ fastest \\
\texttt{infer\_drovFor}          & \texttt{f1:drovFor}             & 2{,}150  & result $\to$ driver + constructor \\
\texttt{infer\_wasTeammate}      & \texttt{f1:wasTeammate}         & 10{,}012 & same race + constructor \\
\texttt{infer\_heldAt}           & \texttt{f1:heldAt}              & 1{,}125  & materialise f1:circuit \\
\texttt{infer\_wonChampionship}  & \texttt{f1:wonChampionship}     & 75       & position=1, MAX(round) \\
\texttt{reify\_wonChampionship}  & \texttt{rdf:Statement} nodes    & 75       & annotate with finalPoints \\
\texttt{infer\_wonCtorChamp.}    & \texttt{f1:wonCtorChampionship} & 17       & same pattern \\
\texttt{classify\_WorldChampion} & \texttt{rdf:type f1:WorldChampion}       & 34  & wonChampionship exists \\
\texttt{classify\_Multichamp.}   & \texttt{rdf:type f1:MultichampionDriver} & 17  & COUNT $\geq$ 2 \\
\texttt{classify\_Veteran}       & \texttt{rdf:type f1:Veteran}             & 86  & COUNT(results) $\geq$ 100 \\
\texttt{classify\_ActiveCircuit} & \texttt{rdf:type f1:ActiveCircuit}       & 24  & race in MAX(year) \\
\texttt{classify\_HistoricCircuit}& \texttt{rdf:type f1:HistoricCircuit}    & 45  & NOT EXISTS in last 10 yrs \\
\texttt{classify\_PodiumResult}  & \texttt{rdf:type f1:PodiumResult}        & 3{,}397 & positionOrder $\leq$ 3 \\
\texttt{classify\_CtorChampion}  & \texttt{rdf:type f1:ConstructorChampion} & 17  & wonCtorChampionship \\
\midrule
\multicolumn{2}{l}{\textbf{Total new triples}} & \textbf{22{,}921} & (verified live in GraphDB) \\
\bottomrule
\end{tabularx}
\end{center}


\subsection{GraphDB Validation}
\label{sec:graphdb-validation}

After loading the integrated RDF file and executing all 21 SPIN rules via
\texttt{python manage.py run\_spin\_rules}, the following triple counts were
verified by SPARQL query against the live GraphDB repository
(\texttt{ws-formula1-owlmax}, OWL-Max ruleset, GraphDB 11.3.1):

\begin{center}
\begin{tabular}{@{}lrl@{}}
\toprule
\textbf{Inferred property / class} & \textbf{Triples} & \textbf{Verification query} \\
\midrule
\texttt{f1:wonRace}           & 1{,}128 & \texttt{COUNT(*) WHERE \{ ?d f1:wonRace ?r \}} \\
\texttt{f1:finishedSecond}    & 1{,}134 & \texttt{COUNT(*) WHERE \{ ?d f1:finishedSecond ?r \}} \\
\texttt{f1:finishedThird}     & 1{,}135 & \texttt{COUNT(*) WHERE \{ ?d f1:finishedThird ?r \}} \\
\texttt{f1:achievedPodium}    & 3{,}395 & \texttt{COUNT(*) WHERE \{ ?d f1:achievedPodium ?r \}} \\
\texttt{f1:startedFromP1}     & 1{,}135 & \texttt{COUNT(*) WHERE \{ ?d f1:startedFromP1 ?r \}} \\
\texttt{f1:convertedP1ToWin}  & 486     & \texttt{COUNT(*) WHERE \{ ?d f1:convertedP1ToWin ?r \}} \\
\texttt{f1:setFastestLap}     & 410     & \texttt{COUNT(*) WHERE \{ ?d f1:setFastestLap ?r \}} \\
\texttt{f1:achievedHatTrick}  & 69      & \texttt{COUNT(*) WHERE \{ ?d f1:achievedHatTrick ?r \}} \\
\texttt{f1:wonChampionship}   & 75      & \texttt{COUNT(*) WHERE \{ ?d f1:wonChampionship ?s \}} \\
\texttt{f1:wasTeammate}       & 10{,}012& \texttt{COUNT(*) WHERE \{ ?d f1:wasTeammate ?d2 \}} \\
\midrule
\texttt{rdf:type f1:WorldChampion} & 34 & \texttt{COUNT(DISTINCT ?d) WHERE \{ ?d rdf:type f1:WorldChampion \}} \\
\texttt{rdf:type f1:Veteran}       & 86 & \texttt{COUNT(DISTINCT ?d) WHERE \{ ?d rdf:type f1:Veteran \}} \\
\texttt{rdf:type f1:PodiumResult}  & 3{,}397 & \texttt{COUNT(*) WHERE \{ ?r rdf:type f1:PodiumResult \}} \\
\midrule
\multicolumn{2}{l}{Total triples in repository} & 6{,}650{,}355 \\
\midrule
\multicolumn{2}{l}{Triples before SPIN (after OWL-Max import)} & 6{,}627{,}434 \\
\multicolumn{2}{l}{\textbf{New triples added by SPIN (21 rules)}} & \textbf{22{,}921} \\
\bottomrule
\end{tabular}
\end{center}

\textbf{Consistency check:} \texttt{f1:wonRace} + \texttt{f1:finishedSecond} + \texttt{f1:finishedThird}
= $1{,}128 + 1{,}134 + 1{,}135 = 3{,}397$, matching \texttt{f1:achievedPodium} (3{,}395)
within a difference of 2, which corresponds to sprint-race results excluded by the
\texttt{FILTER(!CONTAINS(STR(?result), "/sprint-result/"))} guard in the new rules
but not in the original \texttt{infer\_achievedPodium} rule.

The two new rules each completed in under 0.5~seconds (infer\_finishedSecond: 0.19~s,
infer\_finishedThird: 0.16~s). The slowest rule remains \texttt{infer\_wasTeammate}
at 45.8~seconds due to the cross-product join (two drivers per race per constructor).

% =============================================================================
\section{Inference Rules (SPIN)}
\label{sec:spin}
% =============================================================================

While the OWL ontology handles static classification well, the F1 domain requires
patterns that fall outside OWL 2 DL's expressivity. Queries such as ``drivers with
at least 100 race entries'' or ``the driver in position~1 at the final round of a
season'' rely on aggregation and arithmetic, which have no equivalent in Description
Logic. SPIN (SPARQL Inference Notation) bridges this gap by allowing plain
SPARQL~1.1 \texttt{INSERT WHERE} statements to materialise derived triples
directly into the graph.

The specific patterns that required SPIN rules rather than OWL axioms:
\begin{itemize}[leftmargin=1.5em]
  \item \textbf{Aggregation}: \texttt{COUNT}, \texttt{GROUP BY}, \texttt{MAX}
  \item \textbf{Arithmetic comparisons} on data values: \texttt{positionOrder $\leq$ 3}
  \item \textbf{Negation-as-failure}: \texttt{FILTER NOT EXISTS} over derived patterns
  \item \textbf{Cross-entity joins}: same race, same constructor, two distinct drivers
\end{itemize}

All 21 rules are defined in \texttt{championship/spin\_rules.py} as a list of
Python dictionaries, each containing a \texttt{name}, \texttt{description}, and
\texttt{query} (a SPARQL 1.1 \texttt{INSERT WHERE} statement). They are executed
sequentially by the management command \texttt{python manage.py run\_spin\_rules}.

\subsection{Championship Inference}

The most important pair of rules infers the \texttt{f1:wonChampionship} property
by identifying the driver who held position~1 in the \texttt{DriverStanding}
at the \emph{final round} of each season. OWL cannot express this because
``final round'' requires computing the maximum over a GROUP BY aggregate:

\begin{lstlisting}[style=sparql, caption={SPIN rule: infer\_wonChampionship}]
INSERT { ?driver f1:wonChampionship ?season . }
WHERE {
  ?ds f1:driverStandingsId ?anyId ;
      f1:race ?race ;
      f1:driver ?driver ;
      f1:position 1 .
  ?race f1:round ?round ;
        f1:year ?yr .
  ?season rdf:type f1:Season ;   -- type guard: prevents matching Race entities
          f1:year ?yr .
  {
    SELECT ?yr (MAX(?r) AS ?maxRound)
    WHERE { ?rc f1:round ?r ; f1:year ?yr . }
    GROUP BY ?yr
  }
  FILTER(?round = ?maxRound)
  FILTER NOT EXISTS { ?driver f1:wonChampionship ?season }
}
\end{lstlisting}

\begin{infobox}
\textbf{Bug found and fixed:} The initial version lacked the
\texttt{?season rdf:type f1:Season} type guard. Because \texttt{Race} entities
also carry \texttt{f1:year}, the variable \texttt{?season} matched all 17 races
per season in addition to the single \texttt{Season} resource, producing
$17 \times 75 = 1{,}275$ triples instead of 75. Adding the type guard corrected
the count to exactly 75 (one per F1 season from 1950 to 2024).
The same bug affected \texttt{infer\_wonConstructorChampionship} and was fixed
identically.
\end{infobox}

\subsection{Teammate Inference}

The \texttt{infer\_wasTeammate} rule establishes 10{,}012 symmetric pairs between
drivers who shared a constructor in the same race. This is a three-entity cross-join
that OWL cannot express:

\begin{lstlisting}[style=sparql, caption={SPIN rule: infer\_wasTeammate}]
INSERT { ?d1 f1:wasTeammate ?d2 . ?d2 f1:wasTeammate ?d1 . }
WHERE {
  ?r1 f1:resultId ?id1 ; f1:race ?race ; f1:constructor ?ctor ; f1:driver ?d1 .
  ?r2 f1:resultId ?id2 ; f1:race ?race ; f1:constructor ?ctor ; f1:driver ?d2 .
  FILTER(?d1 != ?d2)
  FILTER NOT EXISTS { ?d1 f1:wasTeammate ?d2 }
}
\end{lstlisting}

Note that the rule inserts \emph{both} directions in the same \texttt{INSERT}, rather
than relying on the \texttt{owl:SymmetricProperty} entailment, to ensure full
materialisation in the graph.

\subsection{Classification Rules}

The following three classification rules illustrate the arithmetic and aggregation
patterns that OWL cannot handle:

\begin{lstlisting}[style=sparql, caption={classify\_MultichampionDriver: COUNT aggregation}]
INSERT { ?driver rdf:type f1:MultichampionDriver . }
WHERE {
  SELECT ?driver (COUNT(?season) AS ?titles)
  WHERE { ?driver f1:wonChampionship ?season . }
  GROUP BY ?driver
  HAVING(?titles >= 2)
}
\end{lstlisting}

\begin{lstlisting}[style=sparql, caption={classify\_Veteran: COUNT results $\geq$ 100}]
INSERT { ?driver rdf:type f1:Veteran . }
WHERE {
  SELECT ?driver (COUNT(?r) AS ?entries)
  WHERE { ?r f1:resultId ?id ; f1:driver ?driver . }
  GROUP BY ?driver
  HAVING(?entries >= 100)
}
\end{lstlisting}

\begin{lstlisting}[style=sparql, caption={classify\_HistoricCircuit: MAX(year) arithmetic + negation}]
INSERT { ?circuit rdf:type f1:HistoricCircuit . }
WHERE {
  ?race f1:circuit ?circuit ; f1:year ?yr .
  {
    SELECT (MAX(?y) AS ?maxYear) WHERE { ?r f1:year ?y . }
  }
  FILTER NOT EXISTS {
    ?recentRace f1:circuit ?circuit ;
                f1:year ?recentYr .
    FILTER(xsd:integer(str(?recentYr)) > xsd:integer(str(?maxYear)) - 10)
  }
}
\end{lstlisting}

\subsection{Podium Decomposition Rules}

Two additional rules decompose the existing \texttt{f1:achievedPodium} aggregate
into specific position properties. Combined with the property hierarchy change
(\texttt{f1:wonRace}~$\sqsubseteq$~\texttt{f1:achievedPodium}), the complete podium chain is now:

\begin{lstlisting}[style=turtle, caption={Updated property subPropertyOf chain (verified by HermiT)}]
f1:finishedSecond rdfs:subPropertyOf f1:achievedPodium .
f1:finishedThird  rdfs:subPropertyOf f1:achievedPodium .
f1:wonRace        rdfs:subPropertyOf f1:achievedPodium .   -- changed: was competedIn
f1:achievedPodium rdfs:subPropertyOf f1:competedIn .
\end{lstlisting}

Both rules exclude sprint-result URIs via \texttt{FILTER(!CONTAINS(STR(?result), "/sprint-result/"))},
which the original \texttt{infer\_achievedPodium} does not, explaining the verified
triple-count difference of~2 (Section~\ref{sec:graphdb-validation}).

\textbf{Verified counts (live GraphDB):}
\texttt{f1:finishedSecond} = 1{,}134 triples (0.49~s),
\texttt{f1:finishedThird} = 1{,}135 triples (0.10~s).
These are now used directly in the driver detail page instead of Python-level
filtering over \texttt{f1:PodiumResult}, completing the full inference-in-views audit.

\subsection{Pole Position and Race Achievement Rules}

Four additional rules materialise pole-related and combined race-achievement properties.
These rules were introduced to power the Pole Leaders and Hat Trick pages described in
Section~7.

\begin{description}[leftmargin=2em]
  \item[\texttt{infer\_startedFromP1}]
    Asserts \texttt{?driver f1:startedFromP1 ?race} when a result entity carries
    \texttt{f1:grid = 1} (sprint-result URIs excluded via a \texttt{FILTER}).
    Generates \textbf{1{,}135 triples}.
    Grid position from the result table is used in preference to the qualifying table
    because it has broader historical coverage across all eras.

  \item[\texttt{infer\_convertedP1ToWin}]
    Chains on the outputs of \texttt{infer\_startedFromP1} and \texttt{infer\_wonRace}:
    a driver who has both \texttt{f1:startedFromP1 ?race} and
    \texttt{f1:wonRace ?race} receives \texttt{f1:convertedP1ToWin ?race}.
    Generates \textbf{486 triples}. This is an example of \emph{layered SPIN
    inference} --- one rule consuming the materialised output of two earlier rules.

  \item[\texttt{infer\_setFastestLap}]
    Asserts \texttt{f1:setFastestLap} when \texttt{f1:rank = 1} on a result entity.
    Generates \textbf{410 triples}.

  \item[\texttt{infer\_achievedHatTrick}]
    Requires all three properties to hold simultaneously for the same driver/race pair:
    \texttt{f1:wonRace}, \texttt{f1:startedFromP1}, and \texttt{f1:setFastestLap}.
    Generates \textbf{69 triples} --- the 69 perfect races in F1 history.
\end{description}

\textbf{Rule ordering dependency.}
The \texttt{infer\_convertedP1ToWin} rule must execute after both
\texttt{infer\_startedFromP1} and \texttt{infer\_wonRace}.
\texttt{infer\_achievedHatTrick} depends on all three prior rules.
The management command \texttt{apply\_all\_rules()} enforces this by iterating
rules in declaration order within \texttt{spin\_rules.py}.

\subsection{RDF Reification}

The \texttt{reify\_wonChampionship} rule creates \textbf{75 \texttt{rdf:Statement}
nodes} --- one per F1 World Drivers' Championship from 1950 to 2024. Each node
annotates the existing \texttt{driver wonChampionship season} triple with the
driver's final points total using the custom property \texttt{f1:finalPoints}:

\begin{lstlisting}[style=sparql, caption={SPIN rule: reify\_wonChampionship}]
INSERT {
  _:stmt rdf:type      rdf:Statement ;
         rdf:subject   ?driver ;
         rdf:predicate f1:wonChampionship ;
         rdf:object    ?season ;
         f1:finalPoints ?pts .
}
WHERE {
  ?driver f1:wonChampionship ?season .
  ?season rdf:type f1:Season ; f1:year ?yr .
  ?ds rdf:type f1:DriverStanding ;   -- type guard: prevents matching Result entities
      f1:driver ?driver ;
      f1:race ?race ;
      f1:points ?pts .
  ?race f1:round ?round ; f1:year ?yr .
  {
    SELECT ?yr (MAX(?r) AS ?maxRound)
    WHERE { ?rc f1:round ?r ; f1:year ?yr . }
    GROUP BY ?yr
  }
  FILTER(?round = ?maxRound)
  FILTER NOT EXISTS {
    ?s rdf:type rdf:Statement ;
       rdf:subject ?driver ; rdf:predicate f1:wonChampionship ; rdf:object ?season
  }
}
\end{lstlisting}

\textbf{Why reification?} The triple \texttt{driver wonChampionship season} already
exists as a plain triple. Reification allows attaching metadata \emph{to that triple
itself} --- not to the driver (who has many seasons) or the season (which has one
champion). This is the canonical RDF use case: making statements about statements.
The \texttt{f1:finalPoints} property is declared in the ontology with
\texttt{rdfs:domain rdf:Statement}, ensuring it only appears on reification nodes.

% =============================================================================
\section{New Operations on the Data (SPARQL)}
% =============================================================================

This section presents SPARQL queries that are \textbf{only possible because of
inferred knowledge} --- none can be answered from the raw CSV data directly.

\subsection{World Champions Who Raced as Teammates}

Combines \texttt{f1:wasTeammate} (SPIN-inferred) with \texttt{rdf:type f1:WorldChampion}
(SPIN-inferred). Both properties are absent from the raw data:

\begin{lstlisting}[style=sparql, caption={Champions who shared a constructor in the same race}]
SELECT ?aLabel ?bLabel WHERE {
  ?a rdf:type f1:WorldChampion ; rdfs:label ?aLabel ; f1:wasTeammate ?b .
  ?b rdf:type f1:WorldChampion ; rdfs:label ?bLabel .
  FILTER(?a != ?b)
  FILTER(STR(?aLabel) < STR(?bLabel))
}
ORDER BY ?aLabel
\end{lstlisting}

Result: Alain Prost \& Ayrton Senna, Prost \& Niki Lauda, Prost \& Nigel Mansell
(the legendary McLaren era confirmed by pure graph traversal).

\subsection{Constructors That Fielded the Most World Champions}

Uses \texttt{f1:hadDriver} (OWL \texttt{owl:inverseOf}-inferred from \texttt{f1:drovFor})
combined with \texttt{rdf:type f1:WorldChampion}:

\begin{lstlisting}[style=sparql, caption={Constructor champion headcount via inverseOf-inferred hadDriver}]
SELECT ?ctorLabel (COUNT(DISTINCT ?d) AS ?champions) WHERE {
  ?ctor f1:hadDriver ?d .          -- inverseOf f1:drovFor, not in raw data
  ?d    rdf:type f1:WorldChampion .
  ?ctor rdfs:label ?ctorLabel .
}
GROUP BY ?ctor ?ctorLabel
ORDER BY DESC(?champions)
LIMIT 8
\end{lstlisting}

Result: McLaren (15), Ferrari (15), Williams (11), Team Lotus (10).

\subsection{SPARQL Property Paths via rdfs:subPropertyOf}

The ontology declares \texttt{wonRace $\sqsubseteq$ competedIn} and
\texttt{achievedPodium $\sqsubseteq$ competedIn}. This enables alternation
path queries that would otherwise require a \texttt{UNION}:

\begin{lstlisting}[style=sparql, caption={Circuits where Hamilton won or achieved a podium}]
SELECT DISTINCT ?cLabel WHERE {
  ?h rdfs:label "Lewis Hamilton" .
  ?h (f1:wonRace | f1:achievedPodium) ?race .
  ?race f1:heldAt ?circuit ; rdfs:label ?cLabel .
}
\end{lstlisting}

\subsection{Reification Query: Championship Final Points}

The only way to retrieve the final championship points for each winner is through
the reified \texttt{rdf:Statement} nodes created by the \texttt{reify\_wonChampionship}
SPIN rule:

\begin{lstlisting}[style=sparql, caption={Championship points retrieved from rdf:Statement reification nodes}]
SELECT ?driverLabel ?yr ?pts WHERE {
  ?stmt rdf:type      rdf:Statement ;
        rdf:predicate f1:wonChampionship ;
        rdf:subject   ?driver ;
        rdf:object    ?season ;
        f1:finalPoints ?pts .
  ?driver rdfs:label ?driverLabel .
  ?season f1:year    ?yr .
}
ORDER BY DESC(?yr)
\end{lstlisting}

Sample results: Lewis Hamilton 2020 (347 pts), Max Verstappen 2023 (575 pts),
Michael Schumacher 2004 (148 pts).

\subsection{Veteran World Champions (Class Intersection)}

OWL class intersection of two independently SPIN-inferred classes:

\begin{lstlisting}[style=sparql, caption={Drivers classified as both WorldChampion and Veteran}]
SELECT ?dLabel WHERE {
  ?d rdf:type f1:WorldChampion ;
     rdf:type f1:Veteran ;          -- 100+ race entries, independent SPIN rule
     rdfs:label ?dLabel .
}
ORDER BY ?dLabel
\end{lstlisting}

\subsection{Two-Hop Teammate Graph Traversal}

The \texttt{owl:SymmetricProperty} on \texttt{wasTeammate} enables arbitrary-length
path queries. This finds all drivers who were teammates of Hamilton's own teammates:

\begin{lstlisting}[style=sparql, caption={Teammates of Hamilton's teammates (2-hop traversal)}]
SELECT DISTINCT ?tLabel WHERE {
  ?h rdfs:label "Lewis Hamilton" .
  ?h f1:wasTeammate/f1:wasTeammate ?t .   -- path operator
  ?t rdfs:label ?tLabel .
  FILTER(?t != ?h)
}
\end{lstlisting}

% =============================================================================
\section{Use and Integration of Wikidata and DBpedia}
% =============================================================================

External enrichment is fetched via \textbf{SPARQLWrapper} from the SPARQL endpoints
of Wikidata (\url{https://query.wikidata.org/sparql}) and DBpedia
(\url{https://dbpedia.org/sparql}). All calls are \textbf{deferred to after page
render} --- a JavaScript \texttt{fetch()} call triggers the enrichment after the
page is visible, injecting data into the hero section asynchronously.
This eliminates the 2--8 second SPARQL timeout from the critical page load path.

\subsection{Driver Enrichment (Wikidata)}

Wikidata is queried for F1 drivers using occupation property
\texttt{wdt:P106 wd:Q10843402} (Formula One driver), matched by label:

\begin{lstlisting}[style=sparql, caption={Wikidata query for driver enrichment with new fields}]
SELECT ?item ?birthPlaceLabel ?description ?image
       ?deathDate ?height ?fatherLabel
       (GROUP_CONCAT(DISTINCT ?childLabel; separator="|") AS ?childNames)
WHERE {
  ?item wdt:P106 wd:Q10843402 ;
        rdfs:label ?label .
  OPTIONAL { ?item wdt:P19  ?birthPlace }   -- place of birth
  OPTIONAL { ?item wdt:P18  ?image }        -- portrait photo
  OPTIONAL { ?item wdt:P570 ?deathDate }    -- date of death
  OPTIONAL { ?item wdt:P2048 ?height }      -- height in cm
  OPTIONAL { ?item wdt:P22  ?father .
             ?father rdfs:label ?fatherLabel
             FILTER(LANG(?fatherLabel) = "en") }
  OPTIONAL { ?item wdt:P40  ?child .
             ?child rdfs:label ?childLabel
             FILTER(LANG(?childLabel) = "en") }
  OPTIONAL { ?item schema:description ?description
             FILTER(LANG(?description) = "en") }
  SERVICE wikibase:label { bd:serviceParam wikibase:language "en" . }
  FILTER(LANG(?label) = "en")
  FILTER(LCASE(STR(?label)) = "ayrton senna")
}
GROUP BY ?item ?birthPlaceLabel ?description ?image ?deathDate ?height ?fatherLabel
LIMIT 1
\end{lstlisting}

New fields added in TP2 (beyond TP1's \texttt{birthPlace}):

\begin{center}
\begin{tabular}{@{}lll@{}}
\toprule
\textbf{Wikidata property} & \textbf{Field} & \textbf{Displayed as} \\
\midrule
P570 & \texttt{deathDate} & \texttt{†} suffix on driver name; lifespan chip \\
P2048 & \texttt{height} & Height chip in hero metadata \\
P22 & \texttt{fatherName} & ``Father: Graham Hill'' (F1 dynasties) \\
P40 & \texttt{childNames} & ``Son: Mick Schumacher'' (pipe-separated list) \\
\bottomrule
\end{tabular}
\end{center}

The death date is particularly significant for historical drivers. Senna (1994),
Ratzenberger (1994), Rindt (1970), Clark (1968), and others are now distinguished
visually with the \texttt{†} symbol. The family links reveal well-known F1 dynasties:
Graham/Damon Hill, Gilles/Jacques Villeneuve, Michael/Mick Schumacher.

\textbf{Portrait photos} are loaded from the Wikipedia REST summary endpoint
(\texttt{/api/rest\_v1/page/summary/\{title\}}) via the \texttt{rdfs:seeAlso}
URL stored in the graph. This is preferred over Wikidata P18 (\texttt{wdt:P18})
because the Wikipedia REST API is faster, not rate-limited, and returns
pre-sized thumbnails suitable for direct use in the UI.

\subsection{Constructor Enrichment (Wikidata)}

Constructors are queried as F1 racing teams (\texttt{wdt:P31 wd:Q41298738}):

\begin{lstlisting}[style=sparql, caption={Wikidata query for constructor enrichment}]
SELECT ?item ?foundingDate ?hqLabel ?logo ?website
       ?countryLabel ?founderLabel ?dissolved WHERE {
  ?item wdt:P31 wd:Q41298738 ;
        rdfs:label ?label .
  OPTIONAL { ?item wdt:P571 ?foundingDate }   -- inception date
  OPTIONAL { ?item wdt:P159 ?hq }             -- headquarters
  OPTIONAL { ?item wdt:P154 ?logo }           -- official logo image
  OPTIONAL { ?item wdt:P856 ?website }        -- official website
  OPTIONAL { ?item wdt:P17  ?country }        -- country of origin
  OPTIONAL { ?item wdt:P112 ?founder }        -- founded by
  OPTIONAL { ?item wdt:P576 ?dissolved }      -- dissolved date
  SERVICE wikibase:label { bd:serviceParam wikibase:language "en" . }
  FILTER(LANG(?label) = "en")
  FILTER(LCASE(STR(?label)) = "ferrari")
}
LIMIT 1
\end{lstlisting}

The team logo (\texttt{P154}) replaces the generic Wikipedia thumbnail on the
constructor detail page. The \texttt{dissolved} date displays a ``Dissolved YYYY''
chip for defunct teams such as Tyrrell, Brabham, and Lotus.

\subsection{Circuit Enrichment (DBpedia)}

DBpedia provides track-specific facts not present in the Ergast dataset.
A key finding during development was that F1 circuits in DBpedia are \textbf{not}
classified as \texttt{dbo:RaceTrack} --- they use generic types (\texttt{owl:Thing},
\texttt{geo:SpatialThing}). Label-matching queries with a type constraint
consistently returned zero results.

The solution is to \textbf{construct the DBpedia resource URI directly} from the
circuit name, bypassing the need for any label or type matching:

\begin{lstlisting}[style=python, caption={Direct URI construction for DBpedia circuit lookup (external\_data.py)}]
resource_uri = "http://dbpedia.org/resource/" + circuit_name.replace(" ", "_")
# "Circuit de Monaco"  -->  dbr:Circuit_de_Monaco
\end{lstlisting}

\begin{lstlisting}[style=sparql, caption={DBpedia query using direct resource URI}]
SELECT ?abstract ?length ?turns ?capacity ?opened ?lapRecord WHERE {
  OPTIONAL { <dbr:Circuit_de_Monaco>  dbo:abstract  ?abstract
             FILTER(LANG(?abstract) = "en") }
  OPTIONAL { <dbr:Circuit_de_Monaco>  dbp:length    ?length    }
  OPTIONAL { <dbr:Circuit_de_Monaco>  dbp:turns     ?turns     }
  OPTIONAL { <dbr:Circuit_de_Monaco>  dbp:capacity  ?capacity  }
  OPTIONAL { <dbr:Circuit_de_Monaco>  dbp:opened    ?opened    }
  OPTIONAL { <dbr:Circuit_de_Monaco>  dbp:lapRecord ?lapRecord }
  FILTER(BOUND(?turns) || BOUND(?capacity) || BOUND(?abstract))
}
LIMIT 1
\end{lstlisting}

\begin{center}
\begin{tabular}{@{}llll@{}}
\toprule
\textbf{DBpedia property} & \textbf{Field} & \textbf{Monaco} & \textbf{Silverstone} \\
\midrule
\texttt{dbp:turns}     & corners          & 12       & 6 \\
\texttt{dbp:capacity}  & spectators       & 37{,}000 & 164{,}000 \\
\texttt{dbp:opened}    & year built       & 1929     & 1948 \\
\texttt{dbp:lapRecord} & fastest lap text & \textit{(varies)} & \textit{(varies)} \\
\texttt{dbp:surface}   & track surface    & Asphalt  & Asphalt \\
\texttt{dbo:thumbnail} & circuit photo    & Wikipedia thumbnail & Wikipedia thumbnail \\
\bottomrule
\end{tabular}
\end{center}

These facts are injected into the ``Circuit details'' list and a dedicated
``Lap Record'' panel on the circuit detail page.
The \texttt{dbo:thumbnail} image is used as the fallback hero photo when no
Wikipedia image is available from the REST summary endpoint.

\subsection{Race Enrichment (DBpedia)}

Individual race articles in DBpedia provide event-specific context not available
in the Ergast dataset. The resource URI is constructed from the race label and year:

\begin{lstlisting}[style=python, caption={DBpedia race URI construction (external\_data.py)}]
# "2024 Abu Dhabi Grand Prix" (year=2024) with year prefix stripped:
label_core = race_label.removeprefix(str(year)).strip()   # "Abu Dhabi Grand Prix"
resource_uri = f"http://dbpedia.org/resource/{year}_{label_core.replace(' ', '_')}"
# -> dbr:2024_Abu_Dhabi_Grand_Prix
\end{lstlisting}

\begin{center}
\begin{tabular}{@{}lll@{}}
\toprule
\textbf{DBpedia property} & \textbf{Field} & \textbf{Displayed as} \\
\midrule
\texttt{dbp:weather}    & race-day conditions & Chip in hero meta (e.g.\ ``Sunny'') \\
\texttt{dbp:attendance} & spectator count     & Chip in hero meta (e.g.\ ``452,055 spectators'') \\
\texttt{dbo:thumbnail}  & race/circuit photo  & Right-side image panel in race hero \\
\texttt{dbo:abstract}   & race summary text   & Description paragraph below hero \\
\bottomrule
\end{tabular}
\end{center}

Coverage is best for modern races (2010--2024). Historical races prior to 2000
often lack weather and attendance data but may still have the abstract and thumbnail.
All DBpedia fields degrade gracefully --- if no data is found the race page renders
normally without the enrichment chips.

% =============================================================================
\section{Publication of Semantic Data (RDFa and Microformats)}
% =============================================================================

Every detail page in the application simultaneously embeds two semantic markup
standards. These serve different parsers and different use cases.

\subsection{RDFa 1.1 Lite}

\textbf{RDFa 1.1 Lite} uses the Schema.org vocabulary, declared with the
\texttt{vocab} attribute. It is consumed by Google's Knowledge Graph, search
engine rich results, and the W3C pyRDFa distiller.

\begin{lstlisting}[style=html, caption={RDFa on the driver detail page}]
<article class="driver-hero h-card"
         vocab="https://schema.org/"
         typeof="Person"
         resource="{{ request.build_absolute_uri }}">
  <div class="driver-hero-name p-name" property="name">
    Lewis Hamilton
  </div>
  <span property="nationality">British</span>
  <time property="birthDate" datetime="1985-01-07">b. 1985-01-07</time>
  <link property="sameAs" href="https://en.wikipedia.org/wiki/Lewis_Hamilton" />
</article>
\end{lstlisting}

\begin{center}
\begin{tabular}{@{}lll@{}}
\toprule
\textbf{Page} & \textbf{typeof} & \textbf{Key properties} \\
\midrule
Driver detail & \texttt{Person} & \texttt{name}, \texttt{birthDate}, \texttt{nationality}, \texttt{sameAs} \\
Constructor detail & \texttt{SportsOrganization} & \texttt{name}, \texttt{nationality}, \texttt{sameAs} \\
Circuit detail & \texttt{SportsActivityLocation} & \texttt{name}, \texttt{latitude}, \texttt{longitude}, \texttt{addressLocality} \\
Race detail & \texttt{SportsEvent} & \texttt{name}, \texttt{startDate}, \texttt{sameAs} \\
Season detail & \texttt{Event} & \texttt{name}, \texttt{startDate} \\
Champions list & \texttt{ItemList} & \texttt{itemListElement} (each item: \texttt{Person}) \\
\bottomrule
\end{tabular}
\end{center}

The \texttt{resource} attribute is set to \texttt{\{\{\ request.build\_absolute\_uri\ \}\}},
giving each entity a canonical URI tied to its page URL.

\subsection{Microformats 2}

\textbf{Microformats 2} requires no vocabulary declaration and is parsed by IndieWeb
tools, social readers, and contact importers. It coexists with RDFa on the same element:

\begin{lstlisting}[style=html, caption={Microformats 2 on the circuit detail page (h-card + h-adr)}]
<article class="hero h-card h-adr"
         vocab="https://schema.org/"
         typeof="SportsActivityLocation"
         resource="{{ request.build_absolute_uri }}">
  <span class="p-name" property="name">Circuit de Monaco</span>
  <span class="p-locality" property="addressLocality">Monte Carlo</span>
  <span class="p-country-name" property="addressCountry">Monaco</span>
  <a class="u-url" property="sameAs"
     href="https://en.wikipedia.org/wiki/Circuit_de_Monaco">Wikipedia</a>
</article>
\end{lstlisting}

\begin{center}
\begin{tabular}{@{}lll@{}}
\toprule
\textbf{Class} & \textbf{Standard} & \textbf{Used on} \\
\midrule
\texttt{h-card}         & hCard  & Driver, Constructor, Circuit pages \\
\texttt{h-event}        & hEvent & Race, Season pages \\
\texttt{h-adr}          & hAdr   & Circuit address block \\
\texttt{p-name}         & hCard/hEvent & Entity name \\
\texttt{p-nationality}  & hCard  & Driver/Constructor nationality \\
\texttt{p-country-name} & hAdr   & Circuit country \\
\texttt{p-locality}     & hAdr   & Circuit city \\
\texttt{dt-bday}        & hCard  & Driver date of birth \\
\texttt{dt-start}       & hEvent & Race/Season start date \\
\texttt{u-url}          & hCard/hEvent & Wikipedia link (\texttt{rdfs:seeAlso}) \\
\bottomrule
\end{tabular}
\end{center}

\subsection{Why Both Standards?}

RDFa integrates with the OWL ontology namespace and is consumed by SPARQL-aware tools
and Google's Knowledge Graph. Microformats 2 requires no vocabulary declaration and
is parsed by IndieWeb tools, social readers, and contact importers without configuration.
Embedding both on the same HTML element demonstrates mastery of both standards and
their respective trade-offs.

\subsection{Microformats 2 Parser Demonstration}

The \texttt{/tools/} page includes a live parser powered by \texttt{mf2py} (a Python
Microformats 2 library). It fetches any page of the site and returns the extracted
JSON, demonstrating the full round-trip:

\begin{lstlisting}[style=shell, caption={Calling the live MF2 parser endpoint}]
GET /api/parse-microformats/?url=http://localhost:9000/drivers/1/

# Response (abbreviated):
{
  "items": [{
    "type": ["h-card"],
    "properties": {
      "name": ["Lewis Hamilton"],
      "nationality": ["British"],
      "bday": ["1985-01-07"],
      "url": ["https://en.wikipedia.org/wiki/Lewis_Hamilton"]
    }
  }]
}
\end{lstlisting}

\subsection{Validation of Published Semantic Data}
\label{sec:validation}

\subsubsection{RDFa 1.1 Validation — Step by Step}

\textbf{Tool:} RDFa Play at \url{https://rdfa.info/play/}

\begin{enumerate}[leftmargin=2em]
  \item Start the application: \texttt{bash scripts/run\_webserver.sh}
  \item In your browser, navigate to \texttt{http://localhost:9000/drivers/1/}
        (Lewis Hamilton detail page).
  \item Open browser DevTools \textrightarrow{} right-click the page \textrightarrow{}
        \textit{View Page Source} (Ctrl+U). Select all (Ctrl+A) and copy.
  \item Open \url{https://rdfa.info/play/} and switch to the \textbf{``HTML''} input tab.
  \item Paste the copied HTML and click \textbf{Parse}.
  \item In the right panel, verify the following triples are present:
\end{enumerate}

\begin{lstlisting}[style=turtle, caption={Expected RDFa triples from /drivers/1/ (Hamilton)}]
<http://localhost:9000/drivers/1/>
    a schema:Person ;
    schema:name          "Lewis Hamilton" ;
    schema:nationality   "British" ;
    schema:birthDate     "1985-01-07" ;
    schema:sameAs        <https://en.wikipedia.org/wiki/Lewis_Hamilton> .
\end{lstlisting}

\textbf{Screenshot to capture:} RDFa Play showing the parsed graph panel with the
\texttt{schema:Person} node and at least 4 properties.

\vspace{0.5em}
Repeat for the following pages and verify the expected Schema.org types:

\begin{center}
\begin{tabular}{@{}lll@{}}
\toprule
\textbf{URL} & \textbf{Expected typeof} & \textbf{Key properties} \\
\midrule
\texttt{/drivers/1/}    & \texttt{schema:Person}                  & name, nationality, birthDate, sameAs \\
\texttt{/constructors/6/} & \texttt{schema:SportsOrganization}     & name, nationality, sameAs \\
\texttt{/circuits/6/}   & \texttt{schema:SportsActivityLocation}  & name, addressLocality, latitude, longitude \\
\texttt{/races/18/}     & \texttt{schema:SportsEvent}             & name, location, startDate \\
\texttt{/seasons/2024/} & \texttt{schema:SportsEvent}             & name, startDate \\
\texttt{/}              & \texttt{schema:Dataset}                 & name, description, license \\
\texttt{/champions/}    & \texttt{schema:ItemList}                & itemListElement (Person $\times$ 34) \\
\bottomrule
\end{tabular}
\end{center}

\subsubsection{Alternative: W3C pyRDFa Distiller (automated)}

For machine-readable validation, the W3C pyRDFa distiller can extract triples
from a public URL without copy-pasting HTML. Since the application runs on
localhost, first export the rendered HTML to a file:

\begin{lstlisting}[style=shell, caption={Export driver page HTML and validate with pyRDFa}]
# Export rendered page (server must be running on port 9000)
curl -s http://localhost:9000/drivers/1/ -o /tmp/driver1.html

# Validate locally using pyRDFa Python library
pip install pyRDFa3 2>/dev/null
python3 -c "
import pyRdfa
g = pyRdfa.pyRdfa().graph_from_source('/tmp/driver1.html')
print(g.serialize(format='turtle'))
"
\end{lstlisting}

The output should include the \texttt{schema:Person} subject with the properties
listed above.

\subsubsection{Microformats 2 Validation — Step by Step}

\textbf{Method A: Built-in parser endpoint (recommended)}

\begin{enumerate}[leftmargin=2em]
  \item With the server running on port 9000, open:
\end{enumerate}

\begin{lstlisting}[style=shell, caption={Call the built-in mf2py parser}]
# In a browser or with curl:
curl "http://localhost:9000/api/parse-microformats/?url=http://localhost:9000/drivers/1/"
\end{lstlisting}

\textbf{Expected JSON response:}

\begin{lstlisting}[style=shell, caption={mf2py output for /drivers/1/ (Hamilton)}]
{
  "items": [{
    "type": ["h-card"],
    "properties": {
      "name":        ["Lewis Hamilton"],
      "nationality": ["British"],
      "bday":        ["1985-01-07"],
      "url":         ["https://en.wikipedia.org/wiki/Lewis_Hamilton"]
    }
  }]
}
\end{lstlisting}

\textbf{Screenshot to capture:} Browser showing the JSON response body with
\texttt{"type": ["h-card"]} and the four properties.

\vspace{0.5em}
\textbf{Method B: microformats.io online parser}

\begin{enumerate}[leftmargin=2em]
  \item Open \url{https://microformats.io}
  \item Enter a \textbf{public URL} of the site (requires the server to be
        publicly accessible, e.g.\ via \texttt{ngrok}).
  \item Alternatively, paste the raw HTML into the parser's HTML input.
  \item Verify the parsed JSON output matches the expected structure above.
\end{enumerate}

The following pages each return a different Microformats 2 root class:

\begin{center}
\begin{tabular}{@{}lll@{}}
\toprule
\textbf{URL} & \textbf{MF2 root class} & \textbf{Key properties} \\
\midrule
\texttt{/drivers/\{id\}/}      & \texttt{h-card}  & p-name, p-nationality, dt-bday, u-url \\
\texttt{/constructors/\{id\}/} & \texttt{h-card}  & p-name, p-country-name, u-url \\
\texttt{/circuits/\{id\}/}     & \texttt{h-card h-adr} & p-name, p-locality, p-country-name \\
\texttt{/races/\{id\}/}        & \texttt{h-event} & p-name, p-location, dt-start \\
\texttt{/seasons/\{year\}/}    & \texttt{h-event} & p-name, dt-start \\
\bottomrule
\end{tabular}
\end{center}

\subsubsection{Evidence to Include in Submission}

The following screenshots should be captured from the running application:

\begin{enumerate}[leftmargin=2em]
  \item \textbf{RDFa Play} showing extracted triples for the driver detail page
        (\texttt{/drivers/1/} — schema:Person with 4+ properties visible).
  \item \textbf{RDFa Play} showing the home page (\texttt{/}) parsed as schema:Dataset.
  \item \textbf{Browser JSON response} from \texttt{/api/parse-microformats/?url=.../drivers/1/}
        confirming \texttt{h-card} with \texttt{p-name}, \texttt{p-nationality}, \texttt{dt-bday}.
  \item \textbf{Browser JSON response} from \texttt{/api/parse-microformats/?url=.../circuits/1/}
        confirming \texttt{h-card h-adr} with \texttt{p-locality}.
  \item \textbf{Race detail page} showing the SPIN badge on the ``Inferred'' column header
        and DBpedia weather/attendance chips (if available for that race).
  \item \textbf{Driver detail page} showing the ``Semantic Provenance'' panel open with
        OWL-MAX and SPIN badges visible.
\end{enumerate}

% ── Screenshot placeholders for validation evidence ──────────────────────────
\begin{figure}[H]
\centering
% TODO: take screenshot of https://rdfa.info/play/ with driver page source pasted
\fbox{\parbox{0.85\textwidth}{\centering\vspace{0.9cm}
  \textit{[Screenshot --- RDFa Play showing extracted schema:Person triples for\\
   Hamilton (\texttt{/drivers/1/}). Expected: name, nationality, birthDate, sameAs.\\
   Replace this placeholder before submission.]}
  \vspace{0.9cm}}}
\caption{RDFa Play validation of the driver detail page (\texttt{/drivers/1/}).}
\label{fig:rdfa-play-driver}
\end{figure}

\begin{figure}[H]
\centering
% TODO: open http://localhost:9000/api/parse-microformats/?url=http://localhost:9000/drivers/1/
% and take a screenshot of the JSON response showing "h-card"
\fbox{\parbox{0.85\textwidth}{\centering\vspace{0.9cm}
  \textit{[Screenshot --- browser showing JSON response from\\
   \texttt{/api/parse-microformats/?url=.../drivers/1/}\\
   confirming h-card with p-name, p-nationality, dt-bday.\\
   Replace this placeholder before submission.]}
  \vspace{0.9cm}}}
\caption{Microformats 2 parser endpoint (\texttt{/api/parse-microformats/}) confirming
\texttt{h-card} extraction for the driver detail page.}
\label{fig:mf2-parser}
\end{figure}

\begin{figure}[H]
\centering
% TODO: screenshot of /races/<id>/ showing the Semantic Provenance panel
% and the podium cards with SPIN badges
\fbox{\parbox{0.85\textwidth}{\centering\vspace{0.9cm}
  \textit{[Screenshot --- race detail page (\texttt{/races/1144/}) showing the\\
   Semantic Provenance panel with SPIN badges on\\
   f1:wonRace, f1:finishedSecond, f1:finishedThird.\\
   Replace this placeholder before submission.]}
  \vspace{0.9cm}}}
\caption{Semantic Provenance panel on the race detail page, showing SPIN
inference badges for each podium position.}
\label{fig:race-provenance}
\end{figure}

% =============================================================================
\section{Application Functionalities (UI)}
% =============================================================================

This section describes the main features we added in TP2 that were not present in TP1.

\subsection{Pages Driven by Inferred Knowledge}

Eight pages are \textbf{entirely powered by SPIN-inferred classes and properties}.
None of these pages could exist without the ontology and inference layer.

\begin{center}
\begin{tabular}{@{}lllr@{}}
\toprule
\textbf{URL} & \textbf{Page} & \textbf{SPIN property / class} & \textbf{Count} \\
\midrule
\texttt{/champions/}           & World Champions        & \texttt{f1:WorldChampion}          & 34 \\
\texttt{/champions/multi/}     & Multi-Champions        & \texttt{f1:MultichampionDriver}    & 17 \\
\texttt{/champions/veterans/}  & Veterans               & \texttt{f1:Veteran}                & 86 \\
\texttt{/constructors/champions/} & Constructor Champions & \texttt{f1:ConstructorChampion}  & 17 \\
\texttt{/p1/}                  & Pole Position Leaders  & \texttt{f1:startedFromP1}          & 1{,}135 triples \\
\texttt{/hat-tricks/}          & Hat Trick Hall of Fame & \texttt{f1:achievedHatTrick}       & 69 triples \\
\texttt{/hall-of-fame/}        & Hall of Fame dashboard & all inferred classes               & — \\
\bottomrule
\end{tabular}
\end{center}

\subsubsection{Hall of Fame Dashboard (\texttt{/hall-of-fame/})}

We built a landing page that aggregates all inference-powered achievement categories
in one place. At render time it fires five independent SPARQL queries to retrieve
the all-time record holder in each category (most titles, most poles, most hat tricks,
most race entries, most constructor titles), all pulled directly from the materialised
graph. Six category cards show live counts and link to each sub-page.

\subsubsection{Pole Position Leaders (\texttt{/p1/})}

This page is powered by \texttt{f1:startedFromP1} (1{,}135 triples) and
\texttt{f1:convertedP1ToWin} (486 triples), both materialised by SPIN rules.
The conversion rate (pole-to-win percentage) is computed entirely in SPARQL ---
there is no Python arithmetic involved. This was only possible because
\texttt{infer\_startedFromP1} materialises grid position~1 as a first-class
driver$\to$race property, so the query becomes a simple COUNT rather than a JOIN
through the raw result rows. Both table and card views are available; the
user's preference is stored in \texttt{localStorage}.

\subsubsection{Hat Trick Hall of Fame (\texttt{/hat-tricks/})}

Powered by \texttt{f1:achievedHatTrick} (69 triples) --- the 69 perfect Grands Prix
in F1 history where a single driver took pole, the race win, and the fastest lap.
This property is the deepest inference in the system: it requires three independently
materialised properties (\texttt{wonRace}, \texttt{startedFromP1},
\texttt{setFastestLap}) to hold simultaneously for the same driver/race pair,
demonstrating three-layer SPIN chaining.

\subsubsection{Microformats 2 Parser Demonstration}

The \texttt{/api/parse-microformats/} endpoint accepts any URL from the application
and returns the parsed MF2 JSON using the \texttt{mf2py} library.
This allows validators to confirm that \texttt{h-card}, \texttt{h-event},
\texttt{h-adr}, and related properties are correctly emitted in the rendered HTML.

% ── UI screenshots ────────────────────────────────────────────────────────────
\begin{figure}[H]
\centering
% TODO: screenshot of /hall-of-fame/ showing the 6 category cards with live counts
\fbox{\parbox{0.85\textwidth}{\centering\vspace{0.9cm}
  \textit{[Screenshot --- Hall of Fame dashboard (\texttt{/hall-of-fame/}) showing\\
   the 6 category cards with live counts from SPIN inference.\\
   Replace this placeholder before submission.]}
  \vspace{0.9cm}}}
\caption{Hall of Fame dashboard. All counts (champions, veterans, pole leaders,
hat tricks) are queried live from SPIN-inferred graph properties.}
\label{fig:hall-of-fame}
\end{figure}

\begin{figure}[H]
\centering
% TODO: screenshot of /drivers/ showing the SPIN provenance strip and SPIN badges
% on the Wins and Podiums column headers
\fbox{\parbox{0.85\textwidth}{\centering\vspace{0.9cm}
  \textit{[Screenshot --- Drivers list (\texttt{/drivers/}) showing the SPIN\\
   provenance strip and SPIN badges on the Wins / Podiums column headers.\\
   Replace this placeholder before submission.]}
  \vspace{0.9cm}}}
\caption{Drivers list page showing SPIN provenance strip. Win and podium counts
are derived from \texttt{f1:wonRace}, \texttt{f1:finishedSecond}, and
\texttt{f1:finishedThird}, not from raw \texttt{positionOrder} values.}
\label{fig:drivers-spin}
\end{figure}

\subsection{Hero Section Enrichment}

All detail pages (driver, constructor, circuit) now display:

\begin{itemize}[leftmargin=1.5em]
  \item A \textbf{description paragraph} sourced asynchronously from Wikidata or DBpedia,
        with a small source attribution link.
  \item A \textbf{clickable hero photo} that opens a full-size lightbox using the
        native HTML \texttt{<dialog>} element --- no library required.
  \item For \textbf{drivers}: death date (\texttt{†}) displayed in the name, height chip,
        and family dynasty links (father / children who also raced in F1).
  \item For \textbf{constructors}: Wikidata team logo replaces the Wikipedia thumbnail,
        founder and official website injected into the stats list.
  \item For \textbf{circuits}: DBpedia facts (turns, capacity, year opened, lap record)
        injected into the Circuit details card.
\end{itemize}

\subsection{Navigation Structure}

The navigation bar groups pages into logical domains.
A \textbf{Hall of Fame} dropdown consolidates all inference-powered achievement pages:

\begin{itemize}[leftmargin=1.5em]
  \item \textbf{Drivers}, \textbf{Constructors}, \textbf{Seasons}, \textbf{Races},
        \textbf{Circuits} --- top-level direct links.
  \item \textbf{Hall of Fame} $\triangleright$ Overview / World Champions /
        Multi-Champions / Constructor Champions / Veterans / Pole Leaders / Hat Tricks
\end{itemize}

All achievement pages support both a \textbf{card view} (Wikipedia photos via lazy-loaded
API) and a \textbf{table view}, with the user's preference stored in
\texttt{localStorage} per page key.

\subsection{Card Hover Expansion}

Entity cards (drivers, constructors, circuits, champions) expand on hover using
\texttt{transform: scale(1.08)} with \texttt{transform-origin: top center}.
Because CSS transforms do not affect layout flow, adjacent cards are never pushed
--- the expanded card floats over its neighbours via \texttt{z-index: 10}.

\subsection{Deep Links}

Navigation shortcuts were added throughout the application:
\begin{itemize}[leftmargin=1.5em]
  \item Constructor career chips in the driver detail page link to constructor detail pages.
  \item Constructor names in race results link to constructor detail pages.
  \item Circuit names in the season detail calendar link to circuit detail pages.
  \item Race winner constructors in the season calendar link to constructor detail pages.
\end{itemize}

\subsection{Podium Visibility Control}

The ``all-time wins'' podium at the top of the Drivers, Constructors, and Circuits
list pages is now hidden when browsing pages 2 and beyond. On page~1 it provides
a featured highlight; on subsequent pages it would be confusing context, so it is
suppressed with a \texttt{\{if page == 1\}} condition in the template.

% =============================================================================
\section{Conclusions}
% =============================================================================

TP2 successfully demonstrates the full lifecycle of a semantic web application,
extending TP1's data foundation with a complete knowledge layer. One of the most
significant outcomes is the degree to which inference reduces complexity at query
time: every page in the application that displays win counts, podium breakdowns,
championship records, or pole statistics simply reads pre-materialised graph
properties, with no Python arithmetic or complex runtime joins required.

\subsection{Key Achievements}

\begin{itemize}[leftmargin=1.5em]
  \item \textbf{OWL 2 DL ontology} with 27 classes (20 asserted + 7 SPIN-inferred),
        21 object properties, and 7 disjointness axioms.
        Formally verified in both GraphDB (OWL-Max ruleset) and Protégé
        (HermiT 1.4.3, 1{,}333~ms, zero unsatisfiable classes).
  \item \textbf{21 SPIN rules} covering all classification and inference patterns
        that OWL DL cannot express --- aggregation, arithmetic comparisons, and
        negation-as-failure. All 21 pass with verified real-world counts;
        the rules collectively generated 22{,}921 new triples on top of 6.6M facts.
  \item \textbf{RDF reification} of 75 championship-winning triples annotated with
        the driver's final points total --- a concrete demonstration of making
        statements about statements.
  \item \textbf{Live linked-data enrichment} from Wikidata and DBpedia via
        SPARQLWrapper, plus the Wikipedia REST API, with each source contributing
        information that the others cannot provide.
  \item \textbf{Dual semantic markup} (RDFa 1.1 Lite + Microformats~2) on all detail
        pages, with explicit provenance indicators in the UI showing users exactly
        which inferred property or external source produced each piece of data.
  \item \textbf{Eight new application pages} driven entirely by inferred knowledge,
        including the Hall of Fame dashboard, Pole Position Leaders, and
        Hat Trick Hall of Fame, none of which could exist without the inference layer.
\end{itemize}

\subsection{Challenges and Solutions}

Several non-trivial problems were encountered during development, each of which
required careful analysis of the semantic layer:

\begin{itemize}[leftmargin=1.5em]
  \item \textbf{Type guard missing from SPIN rule:} \texttt{infer\_wonChampionship}
        initially produced 1{,}200 triples instead of 75. The problem was that the
        \texttt{?season} variable also matched \texttt{Race} entities, which also
        carry \texttt{f1:year}. Adding \texttt{?season rdf:type f1:Season} as a
        type guard fixed the count. This underlines the importance of being explicit
        about entity types in SPARQL patterns that traverse heterogeneous data.
  \item \textbf{DBpedia doesn't type F1 circuits as \texttt{dbo:RaceTrack}:}
        Our initial label-matching queries returned zero results. We discovered that
        DBpedia uses generic types for F1 circuits, so we switched to constructing
        the resource URI directly from the circuit name ---
        e.g.\ \texttt{dbr:Circuit\_de\_Monaco} --- which worked across all circuits.
  \item \textbf{Reification entity ambiguity:} Both \texttt{DriverStanding} and
        \texttt{Result} entities share the \texttt{f1:driver} and \texttt{f1:points}
        properties, which caused our reification rule to create duplicate
        \texttt{rdf:Statement} nodes. Adding \texttt{?ds rdf:type f1:DriverStanding}
        as a guard restricted the join to the correct entity type.
  \item \textbf{Django URL ordering conflict:} The route \texttt{constructors/champions/}
        was declared after \texttt{constructors/<str:constructor\_id>/} in
        \texttt{urls.py}, so Django matched ``champions'' as a constructor ID and
        returned a 404-style page. Moving the specific pattern above the wildcard
        fixed it. It is a small but illustrative example of how naming things
        semantically (``champion'') can clash with generic identifier patterns.
\end{itemize}

% =============================================================================
\section{Configuration to Run the Application}
% =============================================================================

\subsection{Prerequisites}

\begin{itemize}[leftmargin=1.5em]
  \item Python 3.12+
  \item \href{https://www.ontotext.com/products/graphdb/download/}{GraphDB Free}
        running on \texttt{http://localhost:7200}, with a repository named
        \textbf{\texttt{ws-formula1-owlmax}} created with the
        \textbf{\texttt{owl-max-optimized}} ruleset
  \item Kaggle Formula 1 CSV files in \texttt{data/raw/} (Ergast-style schema)
\end{itemize}

\subsection{Step-by-Step Setup}

\begin{lstlisting}[style=shell, caption={Complete setup from a fresh clone}]
# 1. Virtual environment and dependencies
python3 -m venv venv && source venv/bin/activate
pip install -r requirements.txt

# 2. Environment configuration
cp .env.example .env
# Edit .env:  GRAPHDB_BASE_URL, GRAPHDB_REPOSITORY, DJANGO_SECRET_KEY

# 3. Create GraphDB repository (owl-max-optimized ruleset)
bash scripts/create_graphdb_repo.sh

# 4. Generate RDF facts from CSV (no rdf:type for base classes)
python3 scripts/csv_to_rdf.py
# Output: data/rdf/formula1.nt  (~798 MB)

# 5. Merge ontology + facts into one file (for GraphDB and Protege)
python3 scripts/merge_integrated.py
# Output: data/rdf/formula1_integrated.nt

# 6. Load into GraphDB with server-side import
bash scripts/load_rdf_to_graphdb.sh

# 7. Run all 21 SPIN inference rules (~1 min, infer_wasTeammate takes ~45 s)
python3 manage.py run_spin_rules
# Expected output: 21/21 rules applied successfully.

# 8. Django setup (SQLite for auth only)
python3 manage.py migrate
python3 manage.py createsuperuser

# 9. Start the development server (port 9000)
bash scripts/run_webserver.sh
\end{lstlisting}

Open \url{http://localhost:9000}. Admin panel: \url{http://localhost:9000/admin-panel/login/}.

\textbf{Re-running after an ontology change:}
If \texttt{formula1\_ontology.ttl} is modified, delete \texttt{formula1\_integrated.nt}
and re-run from step~5. The setup script (\texttt{scripts/setup.sh}) detects the
timestamp difference automatically and skips already-completed steps.

\subsection{Protégé Validation}

To validate the ontology in Protégé 5.6.9 with HermiT:

\begin{enumerate}[leftmargin=2em]
  \item Open Protégé and load \texttt{data/rdf/formula1\_integrated.nt}
        (File $\to$ Open from File, format: N-Triples).
  \item Select Reasoner $\to$ HermiT 1.4.3.
  \item Click Start Reasoner. HermiT completes in under 1 second on schema only
        (0 explicit individuals --- all types are inferred from properties).
  \item Navigate to the Entities tab $\to$ Classes to verify the inferred
        class hierarchy matches Section~\ref{sec:hermit}.
  \item No unsatisfiable classes should appear (red background on any class
        would indicate a modelling error).
\end{enumerate}

\begin{infobox}
\textbf{Note on xsd:date / xsd:gYear:} The ontology intentionally omits
\texttt{rdfs:range} declarations for \texttt{f1:dob}, \texttt{f1:date}, and
\texttt{f1:year} because \texttt{xsd:date} and \texttt{xsd:gYear} are not
in the OWL 2 datatype map (Section~\ref{sec:hermit}). The actual data values
in GraphDB remain typed as \texttt{xsd:date}/\texttt{xsd:gYear} — only the
ontology range assertions are omitted to maintain full OWL 2 DL compliance.
\end{infobox}

\subsection{Re-running Inference}

All SPIN rules are idempotent. To re-run after a data change:

\begin{lstlisting}[style=shell]
python manage.py run_spin_rules
\end{lstlisting}

\subsection{Data Files Summary}

\begin{center}
\begin{tabular}{@{}lll@{}}
\toprule
\textbf{File} & \textbf{Purpose} & \textbf{In git?} \\
\midrule
\texttt{data/rdf/formula1\_ontology.ttl} & OWL 2 DL ontology (hand-authored) & Yes \\
\texttt{data/rdf/formula1.nt}            & Facts only, no rdf:type for base classes & No (generated) \\
\texttt{data/rdf/formula1\_integrated.nt} & Ontology + facts merged (for Protégé) & No (generated) \\
\texttt{championship/spin\_rules.py}     & 21 SPIN rules as Python module & Yes \\
\bottomrule
\end{tabular}
\end{center}

\end{document}
